% ============================================================================
% main.tex — master file. Build: pdflatex main && bibtex main && pdflatex x2
% One file per section (paper/sections/) to manage context across sessions.
% SCOPE RULING (T0, 2026-07-27): pure mathematics only. No physics anywhere.
% Every number herein traces to a checker/certificate: see paper/PLAN.md §3.
% ============================================================================
\documentclass[11pt]{amsart}

\usepackage[T1]{fontenc}
% lmodern: scalable Type-1 T1 text fonts. Required here -- without it the T1
% Computer Modern faces resolve to bitmap fonts and microtype's font expansion
% aborts the build ("auto expansion is only possible with scalable fonts").
\usepackage{lmodern}
\usepackage[utf8]{inputenc}
\usepackage{amsmath,amssymb,amsthm}
\usepackage{booktabs}
\usepackage{array}
\usepackage{microtype}
\usepackage[hidelinks]{hyperref}

% ---------------------------------------------------------------------------
% Theorem environments
% ---------------------------------------------------------------------------
\theoremstyle{plain}
\newtheorem{theorem}{Theorem}[section]
\newtheorem{proposition}[theorem]{Proposition}
\newtheorem{lemma}[theorem]{Lemma}
\newtheorem{corollary}[theorem]{Corollary}
\newtheorem{conjecture}[theorem]{Conjecture}
\theoremstyle{definition}
\newtheorem{definition}[theorem]{Definition}
\newtheorem{remark}[theorem]{Remark}
\newtheorem{observation}[theorem]{Observation}
% Epistemic-status environments: this paper distinguishes, in its theorem
% taxonomy, what is kernel-checked, what is exact-symbolic, and what is
% numerically supported. "Numerical Claim" is deliberately NOT a theorem.
\theoremstyle{definition}
\newtheorem{numclaim}[theorem]{Numerical Claim}

% ---------------------------------------------------------------------------
% Notation macros
% ---------------------------------------------------------------------------
\newcommand{\Q}{\mathbb{Q}}
\newcommand{\Z}{\mathbb{Z}}
\newcommand{\C}{\mathbb{C}}
\newcommand{\U}{U}                       % hyperbolic plane lattice
\newcommand{\thetaop}{\theta}            % theta = z d/dz
\newcommand{\Ltwo}{L_2}
\newcommand{\Lthree}{L_3}
\newcommand{\Symtwo}{\operatorname{Sym}^2}
\newcommand{\GL}{\operatorname{GL}}
\newcommand{\SL}{\operatorname{SL}}
\newcommand{\lat}[1]{\langle #1\rangle}
\newcommand{\NS}{\operatorname{NS}}
\newcommand{\artifact}[1]{\texttt{#1}}   % repo artifact names

\begin{document}

\title[Symmetric squares of Cooper's sporadic operators]{A kernel-checked
symmetric-square structure theorem for Cooper's sporadic Ap\'ery-like
operators, and the monodromy lattice of the $s_7$ family}

\author{Xavier Callens}
% Affiliation intentionally left for T0 to fill in; nothing is fabricated.
\email{callensxavier@gmail.com}

\begin{abstract}
Cooper's three sporadic Ap\'ery-like sequences $s_7$, $s_{10}$, $s_{18}$
satisfy order-3 recurrences belonging to a four-parameter template. We prove
--- uniformly in the template parameters $(a,b,c,d)\in\Q^4$, with the proof
verified by the Lean~4 kernel --- that the associated order-3 operator
$\Lthree$ factors as $P_2\cdot\Symtwo(\Ltwo)$ for an order-2 partner operator
$\Ltwo$ whose coefficients are determined explicitly by the template, and,
independently, that the Almkvist--van Straten symmetric-square criterion
$W\equiv 0$ holds identically on the template. The partner separates the three
sporadic candidates arithmetically: only the $s_7$ partner has an integral
holomorphic solution (OEIS A279619), and we make precise the mechanism --- a
\emph{normalized} integral uniformizer --- underlying that integrality. We
further prove, exactly over $\Q$, that $\Ltwo$ and $\Lthree$ are irreducible
and that $\Lthree$ is the minimal-order annihilator of its holomorphic
solution. Finally, we compute the monodromy of the $s_7$ family numerically to
60 digits, recognize the entries exactly (largest residual
$\approx 10^{-59}$), and derive --- by an exact integer pipeline verified by
two independently written implementations --- that the joint
monodromy-invariant lattice is isometric over $\Z$ to $\U\oplus\lat{14}$, in
exact agreement with the lattice $(M_7)^{\perp}$ of Dolgachev's $M_n$-polarized
K3 framework. We are explicit throughout about epistemic status: which
statements are kernel-checked, which are exact symbolic computation, which are
numerically supported, and which are classical results cited. All
computations are reproducible from the accompanying repositories.
\end{abstract}

\maketitle
\tableofcontents

% ============================================================================
% 01-introduction.tex — DRAFTED IN FULL
% ============================================================================
\section{Introduction}\label{sec:intro}

Ap\'ery's proof of the irrationality of $\zeta(3)$ placed a handful of integer
sequences --- defined by three-term recurrences with polynomial coefficients,
yet integral against all apparent odds --- at the center of a circle of ideas
connecting holonomic differential equations, modular forms, and the geometry
of families of algebraic surfaces. The six \emph{sporadic} Ap\'ery-like
sequences of order 2, catalogued by Zagier, and their order-3 companions have
since been objects of sustained experimental and theoretical attention.
Cooper~\cite{Cooper2012} identified three further sporadic sequences, commonly
denoted $s_7$, $s_{10}$ and $s_{18}$, satisfying order-3 recurrences;
Gorodetsky~\cite{Gorodetsky2023} gave closed forms and placed all three inside
a single four-parameter recurrence template. Almkvist and van
Straten~\cite{AlmkvistvanStraten2021} exhibited explicit projective K3 models
whose periods realize the corresponding order-3 operators.

This paper studies the differential operators attached to Cooper's template
from three angles --- structural, arithmetic, and lattice-theoretic --- with a
methodological commitment that we regard as part of the contribution:
\emph{every claim is labeled by how it is known}. Specifically, each statement
in this paper carries exactly one of the following statuses, and the label is
part of the statement's presentation, not a footnote:

\begin{enumerate}
\item[\textbf{(K)}] \emph{Kernel-checked}: proved in Lean~4 and verified by
  its proof kernel against the pinned \texttt{mathlib}
  library~\cite{deMouraUllrich2021,mathlibCPP2020,mathlib4}, with the axiom
  base disclosed in \S\ref{sec:formalization};
\item[\textbf{(E)}] \emph{Exact}: established by exact symbolic computation
  over $\Q$ (no floating point), reproducible by a named checker script in the
  accompanying repositories, each shipped with negative controls
  (\S\ref{sec:reproducibility});
\item[\textbf{(N)}] \emph{Numerically supported}: derived from
  high-precision numerics followed by exact recognition and exact
  post-verification; stated as a \emph{Numerical Claim}, never as a theorem;
\item[\textbf{(C)}] \emph{Cited}: a classical result of the literature, used
  as stated, with the precise statement and location given.
\end{enumerate}

\subsection{Results}

Write $\thetaop = z\,\frac{d}{dz}$. Cooper's template, in the normalization
of~\cite{Gorodetsky2023}, is the family of order-3 operators
\begin{equation}\label{eq:cooper-template}
  \Lthree(a,b,c,d)
  \;=\;
  \thetaop^3 - z\,(2\thetaop+1)\bigl(a\thetaop^2 + a\thetaop + b\bigr)
  + z^2\,\bigl(c(\thetaop+1)^3 + d(\thetaop+1)\bigr),
  \qquad (a,b,c,d)\in\Q^4,
\end{equation}
with $(a,b,c,d) = (13,4,-27,3)$, $(6,2,-64,4)$, $(14,6,192,-12)$ for $s_7$,
$s_{10}$, $s_{18}$ respectively \cite{Cooper2012,Gorodetsky2023}.

Our first result is a uniform structure theorem
(Theorem~\ref{thm:sym2-template}, status \textbf{(K)}): for \emph{every}
parameter choice, $\Lthree(a,b,c,d)$ is, in a division-free
$\thetaop$-form sense made precise in \S\ref{sec:sym2}, the product
$P_2\cdot\Symtwo(\Ltwo)$ of its own leading $\thetaop$-coefficient with the
symmetric square of an order-2 operator $\Ltwo$ whose coefficients are
\emph{determined} by the template --- the partner is solved out of the
identities, leaving a single nontrivial constraint, which holds identically.
Independently, the Almkvist--van Straten symmetric-square criterion $W\equiv 0$
holds identically on the template in cleared-denominator form
(Theorem~\ref{thm:avs}, status \textbf{(K)}). Both facts are single
kernel-checked proofs covering $s_7$, $s_{10}$, $s_{18}$ --- and any future
instance of the template --- simultaneously by specialization. To our
knowledge this is the first kernel-checked verification of a symmetric-square
structure theorem for Picard--Fuchs-type operators in any proof assistant.

The partner operator makes an arithmetic distinction visible
(\S\ref{sec:integrality}): among Cooper's three sporadic sequences, only the
$s_7$ partner has an integral holomorphic solution --- the sequence
$1, 2, 22, 336, 6006, \dots$ (OEIS A279619~\cite{OEIS}). Non-integrality for
$s_{10}$ and $s_{18}$ is kernel-checked by finite witness (status
\textbf{(K)}); integrality for $s_7$ is a theorem of
O'Brien~\cite{OBrien2016} (status \textbf{(C)}), whose applicability to our
partner is kernel-checked at the level of recurrences. We isolate the
mechanism: the modular coordinate $X_7 = \eta_1^3\eta_7^3/z_7^3$ is a
\emph{normalized} integral uniformizer, and normalization --- not mere
integrality of $\eta$-quotients --- is the load-bearing property (status
\textbf{(E)}).

Third (\S\ref{sec:irreducibility}, status \textbf{(E)}): the $s_7$ and
$s_{10}$ partners $\Ltwo$ are irreducible, by a denominator obstruction ---
every attainable residue sum of a hyperexponential logarithmic derivative lies
in $\tfrac12\Z$, while the exponents of $\Ltwo$ at infinity are
$\{\tfrac13,\tfrac23\}$ (resp.\ $\{\tfrac38,\tfrac58\}$) --- and their
monodromy is not dihedral, because the doubled indicial root at $0$ forces a
logarithmic solution and no nontrivial unipotent element lies in the
normalizer of a maximal torus of $\SL_2$. Consequently
$\Lthree = \Symtwo(\Ltwo)$ is irreducible and is the \emph{minimal-order}
annihilator of its holomorphic solution. All steps are exact over $\Q$.

Finally (\S\ref{sec:lattice}), we compute the monodromy representation of the
$s_7$ partner by numerical analytic continuation at 60-digit working
precision and recognize the entries exactly; the largest recognition residual
is ${\approx}10^{-59}$ against an acceptance gate of $10^{-35}$, and the
recognized matrices pass exact structural post-verification (involutivity,
the cusp relation at infinity, existence and uniqueness up to scale of a
joint invariant form, orbit-lattice closure). This step is numerics-backed
and is reported as Numerical Claim~\ref{nc:monodromy} (status \textbf{(N)}),
not as a theorem. Taking the recognized matrices as input, an exact
integer-arithmetic pipeline then derives the joint monodromy-invariant
lattice in its primitive even scaling and proves --- unconditionally, as a
statement about the exhibited Gram matrix (Proposition~\ref{prop:usplit},
status \textbf{(E)}) --- that it is isometric over $\Z$ to
$\U\oplus\lat{14}$: Gram matrix
$\left[\begin{smallmatrix}0&0&-1\\0&14&0\\-1&0&0\end{smallmatrix}\right]$,
determinant $-14$, signature $(2,1)$, discriminant form $\Z/14$ with
$q=\tfrac1{14}$. The isometry is realized by an explicit $\GL_3(\Z)$ base
change and has been verified by \emph{two independently written} exact
implementations; the identical pipeline applied to the $s_{10}$ family
derives determinant $-20$ and $\U\oplus\lat{20}$, so the computation
discriminates between levels. This matches on the nose the lattice
$(M_7)^{\perp} = \U\oplus\lat{14}$ of Dolgachev's $M_n$-polarized K3
framework~\cite{Dolgachev1996}, in which the Picard--Fuchs operator is a
symmetric square by a theorem of Doran~\cite{Doran2000}; \S\ref{sec:framework}
assembles the evidence that the $s_7$ family is $M_7$-polarized, and states
precisely what remains unproven --- in particular, the identification of the
computed monodromy lattice with the transcendental lattice of the family is
presented conditionally, not as a theorem.

\subsection{What is new}

None of the framework results are ours: symmetric-square structure for
$M_n$-polarized families is Doran's theorem, and $\smash{(M_n)^{\perp} =
\U\oplus\lat{2n}}$ is Dolgachev's computation. The contributions of this
paper are: (i) the uniform, kernel-checked template theorem with an
explicitly determined partner, where the literature treats candidates one at
a time and asserts existence rather than exhibiting the partner from the
parameters; (ii) the exact irreducibility/minimality proofs for these
specific operators; (iii) the sharpened integrality mechanism for $s_7$,
including the checked resolution of an ambiguity in the source's defining
equation; (iv) the explicit monodromy lattice of the $s_7$ family with an
exact, doubly-verified integral splitting $\U\oplus\lat{14}$ --- a
\emph{converse-direction} computation: rather than deducing invariants from a
polarization assumption, we compute the invariants of a concretely given
operator family and match them against the framework, in a situation where,
as Doran notes (\cite{Doran2000}, \S6), the classification of rank-19
lattice polarizations is not complete; and (v) the epistemic architecture
itself --- kernel-checked/exact/numerical/cited labeling with reproducible
artifacts and negative controls --- which we offer as a methodological
contribution to experimental mathematics.

\subsection{Organization}

\S\ref{sec:prelim} fixes notation. \S\ref{sec:operators} presents the two
operators for $s_7$ and their exact Riemann schemes. \S\ref{sec:sym2} states
and proves the template structure theorem. \S\ref{sec:irreducibility} proves
irreducibility and minimality. \S\ref{sec:lattice} presents the monodromy
computation and the integral splitting. \S\ref{sec:integrality} treats the
$s_7$ integrality mechanism. \S\ref{sec:framework} relates the results to the
Dolgachev--Doran framework and states the conditional statements and open
questions. \S\ref{sec:formalization} describes the Lean~4 formalization,
including its complete axiom inventory. \S\ref{sec:reproducibility} indexes
the artifacts needed to reproduce every number in this paper.

% ----------------------------------------------------------------------------
% Generated-by: Fable 5 (Stream 1 paper agent) | Verified-by: claims traced to
% paper/PLAN.md #3 rows 1-20 | Reviewed-by: pending T0 (Xavier)
% ----------------------------------------------------------------------------

% ============================================================================
% 02-preliminaries.tex — DRAFTED IN FULL
% ============================================================================
\section{Preliminaries and notation}\label{sec:prelim}

\subsection{Operators}

Throughout, $z$ is the affine coordinate, $D = \frac{d}{dz}$, and
$\thetaop = zD$ is the Euler operator. We work with linear differential
operators with coefficients in $\Q[z]$ (occasionally $\Q(z)$ when a monic
normalization is needed), acting on formal power series and on multivalued
analytic functions on $\mathbb{P}^1$ minus the singular locus.

An order-2 operator will be written in $\thetaop$-form as
\begin{equation}\label{eq:L2-theta-form}
  \Ltwo \;=\; P_2\,\thetaop^2 + P_1\,\thetaop + P_0,
  \qquad P_i\in\Q[z],
\end{equation}
and an order-3 operator as
$\Lthree = c_3\,\thetaop^3 + c_2\,\thetaop^2 + c_1\,\thetaop + c_0$ with
$c_i \in \Q[z]$. For a monic operator in $D$-form,
$L = D^2 + pD + q$ with $p,q$ rational functions, the \emph{symmetric square}
$\Symtwo(L)$ is the unique monic order-3 operator annihilating all products
$y_iy_j$ of solutions of $L$; explicitly
\begin{equation}\label{eq:sym2-formula}
  \Symtwo(D^2+pD+q) \;=\;
  D^3 + 3p\,D^2 + \bigl(2p^2 + p' + 4q\bigr)D + \bigl(4pq + 2q'\bigr).
\end{equation}
Formula~\eqref{eq:sym2-formula} is validated in the formalization against two
examples derived from first principles (\S\ref{sec:formalization}): the
harmonic oscillator, $\Symtwo(D^2+1) = D^3+4D$ (solution products spanned by
$1, \cos 2t, \sin 2t$), and $\Symtwo(D^2+D) = D^3+3D^2+2D$ (solution products
$1, e^{-t}, e^{-2t}$).

\subsection{Riemann schemes and Fuchsian invariants}

For a Fuchsian operator $L$ of order $n$ on $\mathbb{P}^1$, the
\emph{Riemann scheme} records the local exponents of $L$ at each singular
point. The Fuchs relation states that the sum of all exponents, over all
singular points, equals $\binom{n}{2}(s-2)$ where $s$ is the number of
singular points; for an order-3 operator with four singular points the
required total is $6$. A point of \emph{maximal unipotent monodromy} (MUM) is
one whose exponents all coincide (here: $0,0,0$ at $z=0$), so that the local
monodromy is a single Jordan block. We write $W(L)$ for the Wronskian of a
fundamental system of solutions of $L$.

\subsection{Cooper's template and the three sporadic sequences}

Cooper~\cite{Cooper2012} introduced three sequences $s_7$, $s_{10}$, $s_{18}$
satisfying order-3 Ap\'ery-like recurrences; Gorodetsky~\cite{Gorodetsky2023}
subsumed them in the template~\eqref{eq:cooper-template}, whose
$\thetaop$-form coefficients, after expanding and moving powers of $z$ to the
left (the inner factors have constant coefficients, so no commutator terms
arise), are
\begin{equation}\label{eq:cooper-coeffs}
\begin{aligned}
  c_3 &= 1 - 2a\,z + c\,z^2, &
  c_2 &= -3a\,z + 3c\,z^2, \\
  c_1 &= -(a+2b)\,z + (3c+d)\,z^2, &
  c_0 &= -b\,z + (c+d)\,z^2.
\end{aligned}
\end{equation}
The holomorphic solution of $\Lthree(a,b,c,d)$ at $z=0$ normalized by
$u(0)=1$ has coefficients satisfying Cooper's three-term recurrence; for
$(a,b,c,d)=(13,4,-27,3)$ this is $s_7$ (OEIS A183204: $1, 4, 48, 760,
\dots$), for $(6,2,-64,4)$ it is $s_{10}$ (OEIS A005260, the sums
$\sum_k\binom{n}{k}^4$), and for $(14,6,192,-12)$ it is $s_{18}$
\cite{Cooper2012,Gorodetsky2023}. The parameter values are transcribed from
the literature and are pinned, with source citations attached to the
definitions, in the formal development (\S\ref{sec:formalization}).

\subsection{Lattices}

A \emph{lattice} is a free $\Z$-module of finite rank with a nondegenerate
symmetric bilinear form; it is \emph{even} if $\lat{x,x}\in 2\Z$ for all $x$.
$\U$ denotes the hyperbolic plane
$\left[\begin{smallmatrix}0&1\\1&0\end{smallmatrix}\right]$ and $\lat{m}$ the
rank-1 lattice with Gram matrix $(m)$. The \emph{discriminant group} of a
lattice $T$ is $T^\vee/T$, carrying its discriminant quadratic form
$q\colon T^\vee/T\to\Q/2\Z$. Two lattices are isometric over $\Z$ if their
Gram matrices are conjugate under $\GL_n(\Z)$.

For a K3 surface $X$, $\NS(X)\subset H^2(X,\Z)$ is the N\'eron--Severi
lattice, of rank $\rho(X)$, and the \emph{transcendental lattice} $T(X)$ is
its orthogonal complement, of rank $22-\rho(X)$. Following
Dolgachev~\cite{Dolgachev1996}, for $M$ a fixed even lattice of signature
$(1,t)$, an \emph{$M$-polarized} K3 surface is one with a primitive lattice
embedding $M\hookrightarrow \NS(X)$; for the rank-19 lattices
$M_n = \U\oplus E_8^2\oplus\lat{-2n}$, Dolgachev computes
$(M_n)^{\perp} = \U\oplus\lat{2n}$ (\cite{Dolgachev1996}, \S7) and proves
that the coarse moduli space of $M_n$-polarized K3 surfaces is the Fricke
modular curve $H/\Gamma_0(n)+$ (\cite{Dolgachev1996}, Thm 7.1).

% ----------------------------------------------------------------------------
% Generated-by: Fable 5 (Stream 1 paper agent) | Verified-by: eq (coeffs) vs
% Agora/Sequences/PartnerOperators.lean cooperC3..C0; params vs
% CooperRecurrences.lean; Dolgachev statements vs Stream 2 sources-read brief
% | Reviewed-by: pending T0 (Xavier)
% ----------------------------------------------------------------------------

% ============================================================================
% 03-operators.tex — DRAFTED IN FULL
% All exponent/scheme data from Stream 2 certificates L3_RIEMANN_SCHEME.json
% and L3_IRREDUCIBLE.json (exact symbolic checkers with negative controls).
% ============================================================================
\section{The two operators and their exact Riemann schemes}\label{sec:operators}

We now specialize to $s_7$; the $s_{10}$ data are recorded in parallel where
they illuminate the structure. All results in this section have status
\textbf{(E)}: they are exact computations over $\Q$, reproducible by the
checkers named in \S\ref{sec:reproducibility}.

\subsection{The order-3 operator}

For $(a,b,c,d)=(13,4,-27,3)$, the leading $\thetaop$-coefficient of
$\Lthree = \Lthree(13,4,-27,3)$ is
\[
  c_3 \;=\; 1 - 26z - 27z^2 \;=\; -(z+1)(27z-1),
\]
so the finite singular points are $z=-1$ and $z=\tfrac1{27}$, together with
$z = 0$ and $z=\infty$. (For $s_{10}$: $c_3 = 1-12z-64z^2 = -(4z+1)(16z-1)$,
singular at $-\tfrac14$ and $\tfrac1{16}$.)

\begin{proposition}[Riemann scheme of $\Lthree$; status (E)]
\label{prop:L3-scheme}
The operator $\Lthree$ is Fuchsian with Riemann scheme
\[
\begin{Bmatrix}
  0 & -1 & \tfrac1{27} & \infty \\
  \hline
  0 & 0 & 0 & \tfrac23 \\
  0 & \tfrac12 & \tfrac12 & 1 \\
  0 & 1 & 1 & \tfrac43
\end{Bmatrix}
\]
The exponent sum is $6$, as the Fuchs relation requires for an order-3
operator with four singular points; $z=0$ is a point of maximal unipotent
monodromy; and the Wronskians satisfy $W(\Lthree) = W(\Ltwo)^3$ for the
partner operator $\Ltwo$ of \S\ref{sec:sym2}.
\end{proposition}

For $s_{10}$ the scheme is the same at $0$ and at the two finite points
($-\tfrac14$, $\tfrac1{16}$, exponents $0,\tfrac12,1$), with
$\infty$-exponents $\{\tfrac34, 1, \tfrac54\}$ and Fuchs sum again $6$.

\subsection{The order-2 partner}

The partner operator produced by Theorem~\ref{thm:sym2-template} below is,
for $s_7$, in the $\thetaop$-form~\eqref{eq:L2-theta-form}:
\begin{equation}\label{eq:s7-partner}
  P_2 = 1 - 26z - 27z^2, \qquad
  P_1 = -13z - 27z^2, \qquad
  P_0 = -2z - 6z^2.
\end{equation}

\begin{proposition}[Riemann scheme of $\Ltwo$; status (E)]
\label{prop:L2-scheme}
The operator $\Ltwo$ with coefficients~\eqref{eq:s7-partner} is Fuchsian with
Riemann scheme
\[
\begin{Bmatrix}
  0 & -1 & \tfrac1{27} & \infty \\
  \hline
  0 & 0 & 0 & \tfrac13 \\
  0 & \tfrac12 & \tfrac12 & \tfrac23
\end{Bmatrix}
\]
In particular the local exponent differences are $0$ at the MUM point $z=0$,
$\tfrac12$ at both finite singular points, and $\tfrac13$ at infinity. The
exponents of Proposition~\ref{prop:L3-scheme} are exactly the pairwise sums
of these, as they must be for a symmetric square. (For the $s_{10}$ partner:
$\infty$-exponents $\{\tfrac38, \tfrac58\}$, finite exponents as for $s_7$.)
\end{proposition}

\begin{remark}[Implied Fuchsian signature; status (E) for the signature,
supporting evidence for the identification]\label{rem:signature}
Reading the exponent differences of Proposition~\ref{prop:L2-scheme} as
angles of a hyperbolic triangle group datum gives genus $0$, elliptic points
of orders $(2,2,3)$, one cusp, and hyperbolic area $2\pi\cdot\tfrac23$. This
is precisely the signature of the Fricke group $\Gamma_0(7)+$: $\Gamma_0(7)$
has index $8$ in $\mathrm{PSL}_2(\Z)$ (area $2\pi\cdot\tfrac43$), the Fricke
involution halves the area, fuses the two cusps into one and the two order-3
points into one, and its two fixed points (class numbers
$h(-7)=h(-28)=1$) become the two order-2 elliptic points. The identification
of the modular coordinate with a $\Gamma_0(7)+$ Hauptmodul is taken up in
\S\ref{sec:integrality}. The two finite singular points are thus order-2
elliptic points of the quotient curve --- \emph{not} degeneration points of
some elliptic fibration; no such reading is made anywhere in this paper.
\end{remark}

% ----------------------------------------------------------------------------
% Generated-by: Fable 5 (Stream 1 paper agent) | Verified-by: schemes vs
% L3_RIEMANN_SCHEME.json + L3_IRREDUCIBLE.json (this session); partner coeffs
% vs Agora/Sequences/PartnerOperators.lean s7_P2/P1/P0 | Reviewed-by: pending
% T0 (Xavier)
% ----------------------------------------------------------------------------

% ============================================================================
% 04-sym2.tex — the template structure theorem. Claims #1-#5 of paper/PLAN.md §3.
% All statuses (K) trace to Agora/Sequences/PartnerOperators.lean and
% Agora/Swampland/SymSquareC3b.lean (0 sorry, axiom base disclosed in §9).
% ============================================================================
\section{The template structure theorem}\label{sec:sym2}

This section proves that \emph{every} operator of Cooper's
template~\eqref{eq:cooper-template} is the symmetric square of an order-2
operator whose coefficients are determined by the template parameters, and
that the Almkvist--van Straten self-adjointness criterion holds identically on
the template. Both statements have status \textbf{(K)}: they are verified by
the Lean~4 kernel, uniformly in $(a,b,c,d)\in\Q^4$, so that $s_7$, $s_{10}$,
$s_{18}$ --- and any further instance of the template --- are specializations
of a single proof rather than three separate computations.

\subsection{The symmetric square in \texorpdfstring{$\thetaop$}{theta}-form}

The symmetric-square formula~\eqref{eq:sym2-formula} was stated for the
derivation $D=\frac{d}{dz}$. We use it for $\thetaop = zD$ instead. This is
legitimate and requires no re-derivation:

\begin{lemma}[Change of derivation]\label{lem:theta-sym2}
Let $\delta$ be any derivation of a differential field $K$ of characteristic
$0$ and let $K[\delta]$ be the associated skew polynomial ring, so that
$\delta f = f\delta + \delta(f)$ for $f \in K$. For $L = \delta^2 + \alpha
\delta + \beta$ with $\alpha,\beta \in K$, the monic order-3 operator
annihilating all products of solutions of $L$ is
\begin{equation}\label{eq:sym2-theta}
  \Symtwo(L) \;=\;
  \delta^3 + 3\alpha\,\delta^2
  + \bigl(2\alpha^2 + \delta(\alpha) + 4\beta\bigr)\delta
  + \bigl(4\alpha\beta + 2\delta(\beta)\bigr).
\end{equation}
\end{lemma}

\begin{proof}
The classical derivation of~\eqref{eq:sym2-formula} --- set $u = y_iy_j$,
differentiate four times, and eliminate $y_i',y_j'$ using
$\delta^2 y = -\alpha\,\delta y - \beta y$ --- uses only the commutation rule
$\delta f = f\delta + \delta(f)$ and the field operations. It is therefore
valid verbatim for any derivation, in particular for $\delta = \thetaop$ on
$K = \Q(z)$, where $\thetaop(f) = zf'$.
\end{proof}

Throughout this section $\delta = \thetaop$, and we write $\thetaop(f)$ for
the action on a coefficient (as opposed to $\thetaop$ as an operator symbol).
On polynomials, $\thetaop(f) = z f'$; this is exactly the map \artifact{thetaC}
of the formalization.

\subsection{The partner is determined, not chosen}

Write an order-2 operator in the $\thetaop$-form~\eqref{eq:L2-theta-form},
$\Ltwo = P_2\thetaop^2 + P_1\thetaop + P_0$, and suppose we ask for
\begin{equation}\label{eq:factorization-goal}
  \Lthree(a,b,c,d) \;=\; P_2\cdot\Symtwo\!\left(\tfrac{1}{P_2}\Ltwo\right),
\end{equation}
an identity of operators with coefficients in $\Q(z)$. Comparing
$\thetaop$-coefficients of~\eqref{eq:factorization-goal} against the template
coefficients~\eqref{eq:cooper-coeffs}, and using~\eqref{eq:sym2-theta} with
$\alpha = P_1/P_2$, $\beta = P_0/P_2$, gives four conditions. Read in order,
the first three \emph{solve} for the partner and the fourth is then a
constraint carrying all the content:
\begin{equation}\label{eq:four-conditions}
  \underbrace{c_3 = P_2}_{\text{fixes }P_2}, \qquad
  \underbrace{c_2 = 3P_1}_{\text{fixes }P_1}, \qquad
  \underbrace{c_1 = \thetaop(P_1) + 4P_0}_{\text{fixes }P_0}, \qquad
  \underbrace{c_0 = 2\,\thetaop(P_0)}_{\text{a constraint}}.
\end{equation}
Solving the first three against~\eqref{eq:cooper-coeffs} yields the candidate
partner
\begin{equation}\label{eq:generic-partner}
  P_2 = 1 - 2az + cz^2, \qquad
  P_1 = -az + cz^2, \qquad
  P_0 = -\tfrac{b}{2}z + \tfrac{c+d}{4}z^2 .
\end{equation}
Nothing is free at this point. We emphasise that~\eqref{eq:generic-partner} is
a \emph{closed-form} partner exhibited from the template parameters; no
geometric hypothesis and no candidate-by-candidate search enters.

\begin{remark}[{The partner leaves $\Z[z]$}]\label{rem:partner-not-integral}
The halved and quartered coefficients in~\eqref{eq:generic-partner} are
genuine: the partner lives in $\Q[z]$, not $\Z[z]$, even for integral
template parameters. For $s_7$, $s_{10}$, $s_{18}$ the denominators happen to
clear --- $P_0 = -2z-6z^2$, $-z-15z^2$, $-3z+45z^2$ respectively --- but that
is arithmetic accident per candidate, not a property of the template.
\end{remark}

Two identities make~\eqref{eq:factorization-goal} work, and both hold
identically in $(a,b,c,d)$.

\begin{lemma}[Structural collapse of the $\thetaop^2$-coefficient; status (K)]
\label{lem:c2-collapse}
For every $(a,b,c,d)\in\Q^4$, $\;2c_2 = 3\,\thetaop(c_3)$.
\end{lemma}

\begin{lemma}[The magic collapse; status (K)]\label{lem:magic}
For every $(a,b,c,d)\in\Q^4$, the partner~\eqref{eq:generic-partner}
satisfies $\;\thetaop(P_2) = 2P_1$.
\end{lemma}

\begin{proof}[Proof of Lemmas~\ref{lem:c2-collapse} and~\ref{lem:magic}]
Both are polynomial identities in $\Q[z]$ with parameters: $\thetaop(1-2az+cz^2)
= -2az + 2cz^2 = 2(-az+cz^2)$, and $2(-3az+3cz^2) = 3(-2az+2cz^2)$. In the
formalization they are \artifact{partner\_magic} and
\artifact{cooperC2\_eq\_thetaC\_cooperC3}, each closed by \artifact{ring} after
pushing \artifact{Polynomial.derivative} through the algebraic structure, so
the differentiation itself is kernel-computed rather than transcribed.
\end{proof}

Lemma~\ref{lem:c2-collapse} explains why the second condition
of~\eqref{eq:four-conditions} is not an independent assumption: once $P_2$ is
matched to $c_3$, requiring the $\thetaop^2$-residual to vanish
\emph{forces} $\thetaop(P_2) = 2P_1$. Lemma~\ref{lem:magic} is what makes the
identity division-free, as follows. Let
$D_k$ denote the residual between the $\thetaop^k$-coefficient of
$\frac{1}{c_3}\Lthree$ and that of $\Symtwo(\frac{1}{P_2}\Ltwo)$. Clearing
denominators produces the terms $P_1\thetaop(P_2)$ and $2P_0\thetaop(P_2)$;
substituting $\thetaop(P_2) = 2P_1$ turns them into $2P_1^2$ and $4P_0P_1$,
which cancel against the $-2P_1^2$ and $-4P_1P_0$ already present. Every
surviving term then carries exactly one factor of $P_2$, so
\begin{equation}\label{eq:cleared-residuals}
  P_2^2 D_2 = P_2\,(c_2 - 3P_1),\qquad
  P_2^2 D_1 = P_2\,\bigl(c_1 - \thetaop(P_1) - 4P_0\bigr),\qquad
  P_2^2 D_0 = P_2\,\bigl(c_0 - 2\thetaop(P_0)\bigr),
\end{equation}
and the rational-function identity~\eqref{eq:factorization-goal}, which
a priori needs $P_2^2$ in denominators, reduces to three polynomial
identities in $\Q[z]$.

\subsection{The structure theorem}

\begin{theorem}[Cooper-template structure theorem; status (K)]
\label{thm:sym2-template}
Let $(a,b,c,d)\in\Q^4$ and let $P_2,P_1,P_0 \in \Q[z]$ be as
in~\eqref{eq:generic-partner}. Then all four conditions
of~\eqref{eq:four-conditions} hold identically:
\[
  c_3 = P_2,\qquad c_2 = 3P_1,\qquad
  c_1 = \thetaop(P_1) + 4P_0,\qquad c_0 = 2\,\thetaop(P_0).
\]
Consequently, on the open set where $P_2\neq 0$,
\[
  \Lthree(a,b,c,d) \;=\; P_2 \cdot \Symtwo\!\left(\thetaop^2
  + \tfrac{P_1}{P_2}\thetaop + \tfrac{P_0}{P_2}\right),
\]
so every operator of Cooper's template is the leading $\thetaop$-coefficient
of itself times the symmetric square of an explicitly determined order-2
operator.
\end{theorem}

\begin{proof}
The first identity is a definitional equality
(\artifact{partner\_res3}: the polynomial $1-2az+cz^2$ is literally both
$c_3$ and $P_2$). The second is Lemma~\ref{lem:c2-collapse} combined
with~\eqref{eq:generic-partner} (\artifact{partner\_res2}). For the third,
$\thetaop(P_1) = -az + 2cz^2$ and $4P_0 = -2bz + (c+d)z^2$, whose sum is
$-(a+2b)z + (3c+d)z^2 = c_1$ (\artifact{partner\_res1}). For the fourth,
$2\thetaop(P_0) = -bz + (c+d)z^2 = c_0$ (\artifact{partner\_res0}). Only the
fourth carries content: $P_2,P_1,P_0$ are already pinned by the first three,
so the equation either holds or the template is not a symmetric square. It
holds, identically in $(a,b,c,d)$.

The displayed factorization then follows from~\eqref{eq:cleared-residuals}:
each residual $D_k$ satisfies $P_2^2 D_k = P_2 \cdot 0 = 0$, and $P_2 \neq 0$
in the field $\Q(z)$, so $D_2 = D_1 = D_0 = 0$.

In the formalization the four identities are
\artifact{partner\_res3}, \artifact{partner\_res2}, \artifact{partner\_res1},
\artifact{partner\_res0} of \artifact{Agora/Sequences/PartnerOperators.lean},
each proved once for a universally quantified parameter record. Two of them
require relating the halved and quartered coefficients of $P_0$ to their
cleared forms by hand (the constant-embedding $C$ is opaque to
\artifact{ring}); this is done by the auxiliary lemmas
\artifact{C\_half\_b} and \artifact{C\_quarter\_cd}.
\end{proof}

\begin{remark}[Where the theorem lives, and where it does not]
Theorem~\ref{thm:sym2-template} is an identity in $\Q[z]$ and in
$\Q(z)[\thetaop]$. It asserts nothing about modularity, integrality, or
geometry of the partner --- those are the subjects of
\S\ref{sec:integrality} and \S\ref{sec:framework}, and they separate the
three sporadic candidates sharply even though the structure theorem does not.
Keeping ``has a symmetric-square partner'' distinct from ``has a modular
companion'' is the reason the theorem is stated at template level.
\end{remark}

\subsection{The three sporadic instances}

\begin{corollary}[status (K)]\label{cor:sporadic-partners}
Specializing Theorem~\ref{thm:sym2-template} at the parameter values of
\S\ref{sec:prelim} gives the order-2 partners
\[
\begin{array}{llll}
  s_7: & P_2 = 1-26z-27z^2, & P_1 = -13z-27z^2, & P_0 = -2z-6z^2,\\[2pt]
  s_{10}: & P_2 = 1-12z-64z^2, & P_1 = -6z-64z^2, & P_0 = -z-15z^2,\\[2pt]
  s_{18}: & P_2 = 1-28z+192z^2, & P_1 = -14z+192z^2, & P_0 = -3z+45z^2,
\end{array}
\]
with finite singular points $\{-1,\tfrac1{27}\}$, $\{-\tfrac14,\tfrac1{16}\}$
and $\{\tfrac1{12},\tfrac1{16}\}$ respectively.
\end{corollary}

The $s_7$ and $s_{10}$ partners had been obtained in the program by an
independent computer-algebra extraction before the template argument was
found; that the generic construction reproduces them exactly is a genuine
cross-check of~\eqref{eq:generic-partner} rather than a restatement, and it is
recorded as such in the formalization (\artifact{s7\_P2\_eq} and its
siblings). For $s_{18}$ there was no independent extraction to compare
against: its partner is obtained here with no new algebra and no literature
transcription, which matters because no order-2 companion for $s_{18}$ appears
to have been identified previously.

\begin{remark}[Not an instance of the order-2 sporadic template; status (K)]
The partner operators are \emph{not} instances of the classical order-2
Ap\'ery-like template, whose $\thetaop^1$-coefficient is $\thetaop$ of its
$\thetaop^2$-coefficient. The partners instead satisfy $P_1 = \thetaop(P_2)/2$
(Lemma~\ref{lem:magic}); the two templates agree only when $\thetaop(P_2)=0$.
For $s_7$, matching the $\thetaop^2$-coefficient would force
$\thetaop^1$-coefficient $-26z-54z^2$, whereas the actual $P_1$ is
$-13z-27z^2$. This is recorded as a theorem
(\artifact{s7\_P1\_ne\_thetaC\_P2}, \artifact{s10\_P1\_ne\_thetaC\_P2}) so that
the partner is not re-derived by instantiating the wrong template.
\end{remark}

\subsection{The Almkvist--van Straten criterion on the template}

Independently of Theorem~\ref{thm:sym2-template}, the template satisfies the
Almkvist--van Straten self-adjointness criterion for an order-3 operator to be
a symmetric square~\cite{AlmkvistvanStraten2021}. Writing $\Lthree$ in $D$-form
as $\pi_3 D^3 + \pi_2 D^2 + \pi_1 D + \pi_0$ and setting $a_i = \pi_i/\pi_3$,
the criterion is the vanishing of
\begin{equation}\label{eq:avs-W}
  W \;=\; \tfrac13 a_2'' + \tfrac23 a_2 a_2' + \tfrac{4}{27}a_2^3
        + 2a_0 - \tfrac23 a_1 a_2 - a_1'.
\end{equation}
The $\thetaop\to D$ conversion $\thetaop = zD$, $\thetaop^2 = z^2D^2+zD$,
$\thetaop^3 = z^3D^3+3z^2D^2+zD$ applied to~\eqref{eq:cooper-template} gives
\begin{equation}\label{eq:dform-coeffs}
\begin{aligned}
  \pi_3 &= z^3\bigl(1-2az+cz^2\bigr), &
  \pi_2 &= 3z^2\bigl(1-3az+2cz^2\bigr), \\
  \pi_1 &= z\bigl(1-(6a+2b)z+(7c+d)z^2\bigr), &
  \pi_0 &= z\bigl(-b+(c+d)z\bigr).
\end{aligned}
\end{equation}

\begin{theorem}[$W\equiv 0$ on the template; status (K)]\label{thm:avs}
Define $\mathcal{P} \in \Q[z]$ to be the numerator obtained by substituting
$a_i = \pi_i/\pi_3$ into~\eqref{eq:avs-W} and clearing denominators, i.e.
$\mathcal{P} = 27\,\pi_3^3\,W$, expanded as the six-term polynomial expression
\[
\begin{aligned}
  \mathcal{P} \;=\;{}& 9\bigl(\pi_2''\pi_3^2 - \pi_2\pi_3\pi_3''
     - 2\pi_2'\pi_3\pi_3' + 2\pi_2(\pi_3')^2\bigr)
   + 18\,\pi_2\bigl(\pi_2'\pi_3 - \pi_2\pi_3'\bigr)
   + 4\,\pi_2^3 \\
  &+ 54\,\pi_0\pi_3^2
   - 18\,\pi_1\pi_2\pi_3
   - 27\,\pi_3\bigl(\pi_1'\pi_3 - \pi_1\pi_3'\bigr).
\end{aligned}
\]
Then $\mathcal{P} = 0$ in $\Q[z]$ for every $(a,b,c,d)\in\Q^4$. Since $\pi_3$
is a nonzero polynomial for every parameter choice (its lowest coefficient is
that of $z^3$, equal to $1$), this is equivalent to $W\equiv 0$.
\end{theorem}

\begin{proof}
After expanding~\eqref{eq:dform-coeffs} and computing every derivative
formally, $\mathcal{P}$ is a polynomial in $z$ whose coefficients are
polynomials in $(a,b,c,d)$; all of them vanish. The formalization is
\artifact{P\_cleared\_eq\_zero} in
\artifact{Agora/Swampland/SymSquareC3b.lean}, closed by \artifact{ring} after
\artifact{simp} pushes \artifact{Polynomial.derivative} and the constant
embedding through the algebraic structure. The derivatives are genuine
\artifact{Polynomial.derivative} applications rather than hand-transcribed
closed forms, so a transcription error in a differentiation cannot slip past
\artifact{ring}. The three sporadic cases are one-line specializations
(\artifact{P\_cleared\_s7}, \artifact{P\_cleared\_s10},
\artifact{P\_cleared\_s18}).
\end{proof}

\begin{remark}[Two routes, one conclusion]
Theorems~\ref{thm:sym2-template} and~\ref{thm:avs} are logically independent
verifications of the same structural fact via different mechanisms: the first
exhibits the partner and checks a factorization; the second checks a
differential invariant that does not mention the partner at all. The cleared
form $\mathcal{P}$ was, in the program's workflow, derived twice by
independent computer-algebra routes and required to agree before being
encoded in Lean.
\end{remark}

\begin{remark}[Neither theorem discriminates among candidates]
Both statements hold for \emph{all} $(a,b,c,d)$. Symmetric-square structure is
therefore an entry ticket for the whole template and cannot, by itself,
distinguish $s_7$ from $s_{10}$ or $s_{18}$. What does distinguish them is
arithmetic: the integrality of the partner's holomorphic solution
(\S\ref{sec:integrality}).
\end{remark}

\begin{remark}[Validation of the $\Symtwo$ formula; status (K)]
\label{rem:golden-sym2}
Formula~\eqref{eq:sym2-formula} is not taken on trust in the formalization. It
is validated by two golden tests against examples whose symmetric squares are
known from first principles: $\Symtwo(D^2+1) = D^3+4D$ (solutions
$\sin t,\cos t$; products spanned by $1,\cos 2t,\sin 2t$, annihilated by
$D(D^2+4)$) and $\Symtwo(D^2+D) = D^3+3D^2+2D$ (solutions $1,e^{-t}$; products
$1,e^{-t},e^{-2t}$, annihilated by $D(D+1)(D+2)$). These are
\artifact{symSquare\_golden\_harmonic} and \artifact{symSquare\_golden\_exp} in
\artifact{Agora/SymSquare.lean}; see \S\ref{sec:formalization}.
\end{remark}

% ----------------------------------------------------------------------------
% Generated-by: Opus 5 (Stream 1 paper agent) | Verified-by: every identity and
% coefficient checked against Agora/Sequences/PartnerOperators.lean (partnerP2/
% P1/P0, partner_res0-3, partner_magic, cooperC2_eq_thetaC_cooperC3, s*_P*_eq,
% s7_P1_ne_thetaC_P2) and Agora/Swampland/SymSquareC3b.lean (p3/p2/p1/p0,
% P_cleared, P_cleared_eq_zero, P_cleared_s7/s10/s18); golden tests vs
% Agora/SymSquare.lean §3. Claims #1-#5 of paper/PLAN.md §3.
% | Reviewed-by: pending T0 (Xavier)
% ----------------------------------------------------------------------------

% ============================================================================
% 05-irreducibility.tex — irreducibility and minimality. Claim #13 of PLAN.md §3.
% Exponent data from L3_IRREDUCIBLE.json / L3_RIEMANN_SCHEME.json (exact, ℚ).
% ============================================================================
\section{Irreducibility and minimality}\label{sec:irreducibility}

Everything in this section has status \textbf{(E)}: the arguments are exact
over $\Q$, they involve no floating-point step, and they are reproduced by the
checker \artifact{check\_L3\_irreducible\_minimal.py}, which is shipped with
four negative controls comprising sixteen assertions
(\S\ref{sec:reproducibility}). We treat $s_7$ throughout and record the
$s_{10}$ data in parallel, since the same argument applies verbatim.

We work over $\overline{\Q}(z)$; ``irreducible'' means that the operator
admits no proper right factor of positive order in $\overline{\Q}(z)[D]$,
equivalently that the associated differential module is simple. Since our
operators are Fuchsian --- all singular points regular --- irreducibility of
the differential module is equivalent to irreducibility of the monodromy
representation \cite[Ch.~5--6]{vanderPutSinger2003}, and we pass between the
two freely.

\subsection{No hyperexponential solution}

Recall that a nonzero function $y$ is \emph{hyperexponential} over
$\overline{\Q}(z)$ if $y'/y \in \overline{\Q}(z)$. An order-2 operator is
reducible precisely when it has a hyperexponential solution: a right factor of
order $1$ is $D - u$ with $u \in \overline{\Q}(z)$, and its solution
$\exp\int u$ solves the order-2 operator; conversely a hyperexponential
solution $y$ produces the right factor $D - y'/y$.

\begin{lemma}[Denominator obstruction; status (E)]\label{lem:denominator}
Let $L$ be a Fuchsian operator of order $2$ on $\mathbb{P}^1$ whose finite
singular points and whose point $z=0$ have all local exponents in
$\tfrac12\Z$, and suppose no exponent of $L$ at $\infty$ lies in $\tfrac12\Z$.
Then $L$ has no hyperexponential solution, and hence $L$ is irreducible.
\end{lemma}

\begin{proof}
Suppose $y$ is a hyperexponential solution and put $u = y'/y \in
\overline{\Q}(z)$. Because $L$ is Fuchsian, $y$ has moderate growth at every
point of $\mathbb{P}^1$, so $u$ has at worst simple poles and no polynomial
part; write
\[
  u \;=\; \sum_{p} \frac{r_p}{z - p},
\]
the sum over the finitely many poles. At a singular point $p$ of $L$, the
residue $r_p$ is the local exponent of $y$ at $p$, hence lies in
$\tfrac12\Z$ by hypothesis. At an ordinary point $p$ of $L$, $y$ is
meromorphic near $p$ and $r_p$ is its vanishing order, a (possibly negative)
integer; in fact solutions of a Fuchsian operator are holomorphic at ordinary
points, so $r_p \in \Z_{\geq 0}$. Either way $r_p \in \tfrac12\Z$ for every
$p$, whence
\[
  \sum_p r_p \in \tfrac12\Z .
\]
The exponent of $y$ at $\infty$ is $-\sum_p r_p$, and it must be one of the
exponents of $L$ at $\infty$. But no exponent of $L$ at $\infty$ lies in
$\tfrac12\Z$, a contradiction. Hence no hyperexponential solution exists, and
by the remark above $L$ is irreducible.
\end{proof}

\begin{proposition}[$\Ltwo$ is irreducible; status (E)]\label{prop:L2-irred}
The $s_7$ partner $\Ltwo$ of~\eqref{eq:s7-partner} is irreducible over
$\overline{\Q}(z)$. The same holds for the $s_{10}$ partner.
\end{proposition}

\begin{proof}
By Proposition~\ref{prop:L2-scheme} the exponents of the $s_7$ partner are
$\{0,0\}$ at $z=0$ and $\{0,\tfrac12\}$ at each of $z=-1$ and
$z=\tfrac1{27}$ --- all in $\tfrac12\Z$ --- while at $\infty$ they are
$\{\tfrac13,\tfrac23\}$, neither of which lies in $\tfrac12\Z$.
Lemma~\ref{lem:denominator} applies. For $s_{10}$ the finite data are
identical in shape (exponents $\{0,\tfrac12\}$ at $-\tfrac14$ and
$\tfrac1{16}$, $\{0,0\}$ at $0$) and the exponents at $\infty$ are
$\{\tfrac38,\tfrac58\}$, again outside $\tfrac12\Z$.
\end{proof}

The checker computes the least common denominator $N$ of all attainable
residues at runtime rather than assuming it, obtaining $N = 2$ for both
families, and then verifies that \emph{every} exponent at $\infty$ is
obstructed --- not merely one of them.

\subsection{The monodromy is not dihedral}

Let $V$ denote the rank-2 local system of solutions of $\Ltwo$ on
$\mathbb{P}^1 \setminus \{0,-1,\tfrac1{27},\infty\}$, with monodromy group
$G \subseteq \GL(V)$ and projective image $\overline{G} \subseteq \mathrm{PGL}_2$.
We say the monodromy is \emph{dihedral} if $\overline{G}$ is contained in the
image of the normalizer $N(T)$ of a maximal torus $T \subset \SL_2$.

\begin{lemma}[Nontrivial unipotent at the MUM point; status (E)]
\label{lem:unipotent}
The local monodromy of $\Ltwo$ at $z=0$ is unipotent and $\neq \mathrm{Id}$.
\end{lemma}

\begin{proof}
The indicial equation of $\Ltwo$ at $z=0$ has the double root $0$
(Proposition~\ref{prop:L2-scheme}). By the Frobenius method a repeated
indicial root forces a second solution containing $\log z$: with $A(z)$ the
holomorphic solution normalized by $A(0)=1$, a second solution is
$A(z)\log z + g(z)$ with $g$ holomorphic and $g(0)=0$. Continuing once around
$z=0$ fixes $A$ and adds $2\pi i\,A$ to the second solution, so the local
monodromy is a single nontrivial Jordan block, i.e. unipotent and
$\neq\mathrm{Id}$.
\end{proof}

\begin{lemma}[Not dihedral; status (E)]\label{lem:not-dihedral}
$\overline{G}$ is not contained in the image of the normalizer of a maximal
torus of $\SL_2$.
\end{lemma}

\begin{proof}
Let $g \in N(T)$. Either $g \in T$, in which case $g$ is semisimple, or $g$
lies in the nontrivial coset, in which case $g$ conjugates $T$ by inversion
and is conjugate to an antidiagonal matrix; then $g^2 \in T$ is semisimple. Let
$u$ be a lift of the local monodromy at $z=0$. By
Lemma~\ref{lem:unipotent} $u$ is unipotent and $\neq\mathrm{Id}$, hence not
semisimple; and $u^2$ is again unipotent and $\neq\mathrm{Id}$, hence also not
semisimple. So $u$ lies in neither coset, and $\overline{G}\not\subseteq
N(T)/\{\pm 1\}$.
\end{proof}

\subsection{Irreducibility of \texorpdfstring{$\Lthree$}{L3}}

\begin{lemma}[Symmetric square of a rank-2 system]\label{lem:sym2-irred}
Let $V$ be an irreducible rank-2 local system with monodromy group $G$ and
projective image $\overline{G}$. Then $\Symtwo V$ is reducible if and only if
$\overline{G}$ is contained in the normalizer of a maximal torus.
\end{lemma}

\begin{proof}
In characteristic $0$, $V\otimes V \cong \Symtwo V \oplus \wedge^2 V$ and
$V \otimes V \cong \mathrm{End}(V)\otimes \det V$; comparing ranks and the
trace decomposition $\mathrm{End}(V) = \mathrm{End}_0(V)\oplus \mathrm{Id}$
gives
\[
  \Symtwo V \;\cong\; \mathrm{End}_0(V) \otimes \det V ,
\]
so $\Symtwo V$ is reducible if and only if $\overline{G}$ acts reducibly on
$\mathfrak{sl}(V) = \mathrm{End}_0(V)$ by conjugation.

Suppose $\mathfrak{sl}(V)$ has a $\overline{G}$-stable line, spanned by $X\neq
0$. If $X$ is nilpotent, $\ker X \subset V$ is a line and is
$\overline{G}$-stable, contradicting irreducibility of $V$. Hence $X$ is
semisimple with two distinct eigenvalues, and $\overline{G}$ permutes its two
eigenlines in $V$; that is, $\overline{G}$ normalizes the maximal torus
determined by that eigenline pair. If instead $\mathfrak{sl}(V)$ has a stable
plane, its annihilator under the (nondegenerate, invariant) Killing form is a
stable line and the previous case applies. The converse is immediate: if
$\overline{G}$ normalizes $T$ with eigenlines $\langle e_1\rangle, \langle
e_2\rangle$, then the line spanned by $e_1e_2$ in $\Symtwo V$ is stable.
\end{proof}

\begin{theorem}[$\Lthree$ is irreducible; status (E)]\label{thm:L3-irred}
The operator $\Lthree = \Lthree(13,4,-27,3)$ is irreducible over
$\overline{\Q}(z)$. The same holds for $\Lthree(6,2,-64,4)$.
\end{theorem}

\begin{proof}
By Theorem~\ref{thm:sym2-template}, $\Lthree$ is $P_2\cdot\Symtwo(\Ltwo/P_2)$,
so its solution space is the space of products of solutions of $\Ltwo$, i.e.
the local system $\Symtwo V$. $V$ is irreducible by
Proposition~\ref{prop:L2-irred}, and $\overline{G}$ is not contained in the
normalizer of a maximal torus by Lemma~\ref{lem:not-dihedral}. By
Lemma~\ref{lem:sym2-irred}, $\Symtwo V$ is irreducible, hence so is $\Lthree$.
\end{proof}

\begin{corollary}[Minimality; status (E)]\label{cor:minimal}
$\Lthree$ is the minimal-order annihilator in $\overline{\Q}(z)[D]$ of each of
its nonzero solutions; in particular of the holomorphic solution
$u(z) = \sum_{n\geq 0}s_7(n)z^n$ at $z=0$. Consequently the differential
module generated by $u$ has rank exactly $3$.
\end{corollary}

\begin{proof}
Let $y\neq 0$ solve $\Lthree y = 0$ and suppose $M$ has order $m$ with
$1 \leq m < 3$ and $My = 0$. Then $\operatorname{gcrd}(\Lthree,M)$
annihilates $y$, so it
has order $\geq 1$; and it is a right factor of $\Lthree$ of order
$\leq m < 3$. This is a proper right factor, contradicting
Theorem~\ref{thm:L3-irred}. Hence no annihilator of order $<3$ exists, and
$\Lthree$ itself annihilates $y$.
\end{proof}

\begin{remark}[Why minimality is the load-bearing statement]
\label{rem:minimality-load}
Corollary~\ref{cor:minimal} is what pins the rank of the differential module
generated by the holomorphic solution at $3$ rather than at $1$ or $2$. Every
rank-based statement downstream --- in particular the conditional rank
statement of \S\ref{sec:framework} --- rests on it. Had $\Lthree$ factored,
that rank, and everything computed from it, would move. This is why the
argument is given in full here rather than cited.
\end{remark}

\begin{remark}[Negative controls; status (E)]
The checker reproducing this section is deliberately fallible. Its controls
feed it (i) an operator with a hyperexponential solution by construction,
$\thetaop^2-2\thetaop+1$, which annihilates $y = z$ so that $y'/y = 1/z$ is
rational --- the irreducibility step must and does reject it; (ii) an operator
with distinct indicial roots at $0$, so no logarithm and no unipotent --- the
non-dihedral step must and does reject it; (iii) a cross-check against the
independently computed Riemann scheme of \S\ref{sec:operators}, requiring that
the pairwise sums of the $\Ltwo$ exponents reproduce the $\Lthree$ exponents at
all four singular points for both families; and (iv) a positive control that
the real operators still pass. All sixteen assertions behave as required.
\end{remark}

% ----------------------------------------------------------------------------
% Generated-by: Opus 5 (Stream 1 paper agent) | Verified-by: exponent data and
% step structure checked against Stream 2 data/certificates/L3_IRREDUCIBLE.json
% and L3_RIEMANN_SCHEME.json; control inventory checked by running
% checkers/test_L3_irreducible_minimal_controls.py (16 assertions, all ok).
% Claim #13 of paper/PLAN.md §3. | Reviewed-by: pending T0 (Xavier)
% ----------------------------------------------------------------------------

% CONDITIONAL SECTION — included under Option A of paper/PLAN.md §1.3 (scope
% ruling still open for T0). Under Option B this file is deleted wholesale,
% together with 08-dolgachev-doran.tex; no other section depends on it except
% via the labels sec:lattice, nc:monodromy, prop:usplit.
% ============================================================================
% 06-lattice.tex — monodromy and the integral lattice. Claims #14-#17 of
% paper/PLAN.md §3. Every number traces to C2_cooper_s7_v4.json (SHA256
% 036cd895db892aee9802514ec18668535f8896cee53cf6123533af9844387c3e).
% ============================================================================
\section{The monodromy lattice of the \texorpdfstring{$s_7$}{s7} family}
\label{sec:lattice}

\subsection*{Disclosure}

\emph{The monodromy matrices used in this section are obtained numerically.}
They are computed by analytic continuation at 60-digit working precision and
then recognized as exact algebraic numbers. That recognition step is not a
proof, and nothing derived from it is stated as a theorem here: the recognized
matrices appear as Numerical Claim~\ref{nc:monodromy}, status \textbf{(N)}.
Everything downstream of them --- the invariant form, the orbit lattice, the
Gram matrix, and the integral splitting --- is exact integer and rational
arithmetic, and is stated as a proposition \emph{about the recognized
matrices}, i.e. conditionally on Numerical Claim~\ref{nc:monodromy}. The one
statement in this section that is unconditional, because it is a finite
verifiable assertion about an explicitly displayed integer matrix, is
Proposition~\ref{prop:usplit}. We are equally explicit in
\S\ref{sec:framework} about what is \emph{not} claimed: nothing here
identifies the computed lattice with any invariant of a surface.

\subsection{Numerical continuation}

Work in the Frobenius basis of $\Ltwo$ at the MUM point $z=0$: the
holomorphic solution $A(z)$ with $A(0)=1$, and
\[
  \widehat{B}(z) \;=\; \frac{A(z)\log z + g(z)}{2\pi i},
  \qquad g \text{ holomorphic},\ g(0)=0,
\]
whose existence is Lemma~\ref{lem:unipotent}. In this basis the local
monodromy at $z=0$ is $\left[\begin{smallmatrix}1&1\\0&1\end{smallmatrix}\right]$
by construction.

The finite singular loci are computed at runtime as the roots of the leading
$\thetaop$-coefficient of $\Ltwo$, giving $\{-1,\tfrac1{27}\}$ for $s_7$; they
are not typed in. Continuation is by Taylor stepping at 60 decimal digits of
working precision, series order $140$, step ratio $0.35$, with the series tail
bounded by $10^{-45}$ at every step and the bound asserted rather than
assumed. Loops are taken based at $z_0 = \rho/4$, where $\rho$ is the distance
from $0$ to the nearest singular point.

\begin{remark}[Machinery control]\label{rem:machinery-control}
Before any unknown loop is computed, the numerical machinery is required to
reproduce a matrix known analytically: the continued loop around $z=0$ must
return $\left[\begin{smallmatrix}1&1\\0&1\end{smallmatrix}\right]$. It does, to
$3.17\times 10^{-61}$. A pipeline that cannot recover a known answer is not
used to produce unknown ones.
\end{remark}

\begin{numclaim}[Elliptic monodromies of the $s_7$ partner; status (N)]
\label{nc:monodromy}
In the Frobenius basis above, the monodromy matrices of $\Ltwo$ around the two
finite singular points are
\[
  M(\tfrac1{27}) \;=\; \frac{i}{\sqrt7}
  \begin{pmatrix} 0 & 1 \\ -7 & 0\end{pmatrix},
  \qquad
  M(-1) \;=\; \frac{i}{\sqrt7}
  \begin{pmatrix} 7 & 4 \\ -14 & -7\end{pmatrix} .
\]
\end{numclaim}

The evidence, and the limits of that evidence, are as follows. The two
numerically continued matrices satisfy $\det M = -1$ and $M^2 = \mathrm{Id}$ to
about $10^{-59}$, as the local exponents $\{0,\tfrac12\}$ of
Proposition~\ref{prop:L2-scheme} require. The quantity $\sqrt7$ is not put in
by hand: it emerges from the continuation. What is actually recognized is not
the $2\times 2$ matrices themselves but their symmetric squares in the flag
basis $(A^2, A\widehat{B}, \widehat{B}^2)$, which are rational; recognition is
performed with a loud acceptance gate of $10^{-35}$, the largest observed
residual is ${\approx}10^{-59}$, and the denominators appearing in the
recognized entries come out as $\{1,7\}$. The recognized matrices are then
subjected to exact post-verification: each is an involution; the product of
the three loops matches the monodromy at $\infty$ predicted by the
independently computed exponents $\{\tfrac13,\tfrac23\}$ there (trace $0$,
order $3$); the joint invariant symmetric bilinear form exists and is unique
up to scale; and the orbit lattice below closes under all generators and their
inverses. A conspiracy of numerical errors surviving all of these exact gates
is not credible. It is nevertheless not a proof, which is why this is a
Numerical Claim.

\subsection{The exact lattice pipeline}

From here on, all arithmetic is exact --- integers and rationals only.

\begin{proposition}[The joint monodromy-invariant lattice; status (E), given
Numerical Claim~\ref{nc:monodromy}]\label{prop:lattice}
Let $N(p) = \Symtwo M(p)$ denote the recognized rational $3\times3$ matrices
and let $\Gamma$ be the group they generate together with the cusp matrix
$N(0) = \left[\begin{smallmatrix}1&1&1\\0&1&2\\0&0&1\end{smallmatrix}\right]$.
Then:
\begin{enumerate}
\item the space of $\Gamma$-invariant symmetric bilinear forms on $\Q^3$ is
  one-dimensional;
\item the $\Gamma$-orbit of the cusp-invariant isotropic vector $f$ spans a
  lattice that stabilizes after finitely many steps and is closed under all
  generators and inverses;
\item in its primitive even rescaling, the restriction of the invariant form
  to that lattice has Gram matrix
  \[
    G \;=\;
    \begin{pmatrix} 0&0&-1\\ 0&14&0\\ -1&0&0 \end{pmatrix},
  \]
  with $\det G = -14$, signature $(2,1)$, discriminant group $\Z/14$ and
  discriminant form (in the sense of~\cite{Nikulin1979}) taking the value
  $q = \tfrac{1}{14} \bmod 2\Z$ on a generator;
\item the divisibility of $(N(0)-\mathrm{Id})^2$ on the lattice yields
  $2n = 14$;
\item the lattice admits no proper even $\Gamma$-invariant overlattice.
\end{enumerate}
\end{proposition}

Items (i)--(v) are finite exact computations on the recognized matrices; they
are performed by \artifact{check\_U1\_lattice.py}, whose structural conditions
are asserted (and therefore can fail) rather than reported. The sign of the
form is fixed by an independent positivity re-derivation of the framework
period shape, not chosen.

\subsection{The integral splitting}

The next statement is the one place in this section where nothing is
conditional. It is a finite assertion about the explicit integer matrix $G$
displayed above, checkable by hand in a minute, and we give the witness so
that the reader can do so.

\begin{proposition}[$G \cong \U\oplus\lat{14}$ over $\Z$; status (E),
unconditional]\label{prop:usplit}
Let $G$ be the integer matrix of Proposition~\ref{prop:lattice}(iii) and let
\[
  P \;=\;
  \begin{pmatrix} 1&0&0\\ 0&0&-1\\ 0&-1&0\end{pmatrix},
  \qquad \det P = -1 .
\]
Then $P \in \GL_3(\Z)$ and
\[
  P^{\mathsf T} G P \;=\;
  \begin{pmatrix} 0&1&0\\ 1&0&0\\ 0&0&14 \end{pmatrix}
  \;=\; \U \oplus \lat{14}.
\]
In particular the lattice with Gram matrix $G$ is isometric over $\Z$ to
$\U\oplus\lat{14}$.
\end{proposition}

\begin{proof}
Write $p_1,p_2,p_3$ for the columns of $P$ and $\langle\cdot,\cdot\rangle$ for
the form with Gram matrix $G$ in the standard basis $e_1,e_2,e_3$. Then
$p_1 = e_1$, $p_2 = -e_3$, $p_3 = -e_2$, and directly from $G$:
$\langle p_1,p_1\rangle = G_{11} = 0$;
$\langle p_1,p_2\rangle = -G_{13} = 1$;
$\langle p_1,p_3\rangle = -G_{12} = 0$;
$\langle p_2,p_2\rangle = G_{33} = 0$;
$\langle p_2,p_3\rangle = G_{32} = 0$;
$\langle p_3,p_3\rangle = G_{22} = 14$.
So $\{p_1,p_2\}$ spans a hyperbolic plane $\U$, orthogonal to $\langle
p_3\rangle \cong \lat{14}$. Expanding $\det P$ along the first column gives
$\det\left[\begin{smallmatrix}0&-1\\-1&0\end{smallmatrix}\right] = -1$, so $P$
is unimodular and the base change is an isometry of lattices over $\Z$.
\end{proof}

\begin{remark}[Two independent verifications]\label{rem:two-verifications}
Proposition~\ref{prop:usplit} has been verified by two implementations written
independently of one another and sharing no code.
\begin{enumerate}
\item The pipeline of \artifact{check\_U1\_lattice.py} constructs a base change
  of determinant $+1$ during its stage 3 and checks $S^{\mathsf T}GS =
  \U\oplus\lat{14}$ exactly. It serializes $\det S$ and $S^{\mathsf T}GS$ to
  its certificate but not $S$ itself.
\item A from-scratch implementation, \artifact{check\_U1\_splitting\_independent.py},
  loads the certificate at runtime, reads only $G$ and the integer $d$ (never
  hardcoding $14$ outside a labeled control), assembles the target
  $\U\oplus\lat{d}$ at runtime, and re-derives its own witness from $G$ alone
  --- locating a primitive isotropic vector of unimodular pairing, completing
  it to a hyperbolic pair by an extended-gcd combination, and taking the
  orthogonal-complement generator. All arithmetic is exact Python integers. It
  returns the matrix $P$ displayed in Proposition~\ref{prop:usplit}, of
  determinant $-1$.
\end{enumerate}
The two witnesses differ (determinants $+1$ and $-1$); that they differ is
expected, since any element of $\GL_3(\Z)$ realizing the isometry is an equally
valid proof, and it confirms that the second implementation did not reproduce
the first. The second is shipped with five negative controls: perturbing one
entry of $G$ must break agreement with the certificate's own determinant, and
does; corrupting the witness must break $P^{\mathsf T}GP = \U\oplus\lat{d}$,
and does; doubling a column of the witness must be rejected at the
$\GL_3(\Z)$ gate, and is; testing the genuine $s_7$ Gram against the target
$\U\oplus\lat{20}$ must fail, and does; and the unmodified certificate must
still pass after all scrambling hooks have been exercised, and does.
\end{remark}

\subsection{The \texorpdfstring{$s_{10}$}{s10} control}

\begin{numclaim}[Level discrimination; status (N)+(E)]\label{nc:s10-control}
The identical pipeline --- same code, same tolerances, same exact
post-verification --- applied to the $s_{10}$ family derives determinant $-20$,
$2n = 20$, and an integral splitting $\U\oplus\lat{20}$.
\end{numclaim}

This is the control that matters most. A pipeline that produced
$\U\oplus\lat{14}$ from any input, or that had $14$ wired into it anywhere,
would be worthless; running it unchanged on a different member of the same
template and obtaining a different, internally consistent answer at a
different level shows that the computation is sensitive to the operator it is
given. The certificate records the $s_{10}$ values as computed-versus-computed
inequalities against the $s_7$ values, with no expected result typed in.
Further controls in the same suite corrupt the recognized matrices in three
distinct ways: a raw entry perturbation fails at the involution gate; a
conjugated involution fails at the infinity-product gate; and a corruption fed
directly to the invariant-form step causes the joint invariant form to vanish
(solution space of dimension $0$) rather than to be silently replaced by a
different form.

\begin{remark}[What has and has not been established]
\label{rem:lattice-scope}
Established, conditionally on Numerical Claim~\ref{nc:monodromy}: the joint
monodromy-invariant lattice of the $s_7$ family, in its primitive even
scaling, is isometric over $\Z$ to $\U\oplus\lat{14}$, with the discriminant
data of Proposition~\ref{prop:lattice}. Established unconditionally:
Proposition~\ref{prop:usplit}, a statement about the displayed matrix $G$.
\emph{Not} established anywhere in this section: that this lattice is the
transcendental lattice of anything. That identification is a separate question,
taken up --- and left conditional --- in \S\ref{sec:framework}.
\end{remark}

% ----------------------------------------------------------------------------
% Generated-by: Opus 5 (Stream 1 paper agent) | Verified-by: all numbers read
% from Stream 2 data/certificates/C2_cooper_s7_v4.json (derived block, controls
% block, how block) and briefs/STREAM2_U1_EXECUTION_2026_07_27.md; witness P and
% control inventory from this repo's briefs/STREAM1_U1_INDEPENDENT_VERIFICATION_
% 2026_07_27.md and checkers/check_U1_splitting_independent.py; P^T G P verified
% by hand in the proof. Claims #14-#17 of paper/PLAN.md §3.
% | Reviewed-by: pending T0 (Xavier)
% ----------------------------------------------------------------------------

% ============================================================================
% 07-integrality.tex — the arithmetic separation. Claims #7-#10 of PLAN.md §3.
% Attribution discipline: integrality of the s7 partner is O'Brien 2016 Thm 6.2
% (CITED), used in the formalization as the single registered literature axiom;
% the recurrence-matching step is kernel-checked. Nothing here re-proves it.
% ============================================================================
\section{The arithmetic separation of the three candidates}
\label{sec:integrality}

Theorem~\ref{thm:sym2-template} is blind to the choice of parameters. This
section exhibits the invariant that is not: the integrality of the holomorphic
solution of the order-2 partner. Among Cooper's three sporadic sequences, only
$s_7$ has an integral partner, and we make the mechanism behind that
integrality precise.

\subsection{The partner recurrence}

Extracting the coefficient of $z^{M}$ from $\Ltwo f = 0$ for
$f = \sum_{n\geq0}a_nz^n$ with $a_0 = 1$, using the partner
coefficients~\eqref{eq:generic-partner}, gives $a_1 = \tfrac{b}{2}$ and, for
$k\geq0$,
\begin{equation}\label{eq:partner-recurrence}
  (k+2)^2 a_{k+2}
  \;=\; a_{k+1}\Bigl(2a(k+1)^2 + a(k+1) + \tfrac{b}{2}\Bigr)
      - a_k\Bigl(ck^2 + ck + \tfrac{c+d}{4}\Bigr).
\end{equation}
The leading coefficient $(k+2)^2$ is the crux: the sequence is
$\Q$-valued by construction, and integrality --- when it holds --- is an
Ap\'ery-style theorem, not a consequence of the shape of the recursion.

There is, however, a floor. Let $s\colon\mathbb{N}\to\Z$ be \emph{any} integer
sequence with $s(0)=1$ and let $g$ be the coefficient sequence of the formal
square root of $\sum_n s(n)z^n$, defined by $g_0=1$ and
$g_n = \bigl(s(n) - \sum_{k=1}^{n-1}g_kg_{n-k}\bigr)/2$.

\begin{proposition}[Dyadic baseline; status (K)]\label{prop:dyadic}
For every integer sequence $s$ with $s(0)=1$, every $g_n$ lies in
$\Z[\tfrac12]$.
\end{proposition}

\begin{proof}
Induction on $n$: $\Z[\tfrac12]$ is a ring containing $\Z$, so the sum of
products of earlier terms lies in it, and it is closed under division by $2$.
Formalized as \artifact{sqrtSeq\_dyadic} in
\artifact{Agora/Sequences/FormalSqrt.lean}; the accompanying control
\artifact{dyadic\_baseline\_control} checks that $\tfrac13\notin\Z[\tfrac12]$,
so the statement is not vacuous.
\end{proof}

Proposition~\ref{prop:dyadic} concerns the formal square root $g$; the object
of interest is the recurrence-defined partner $a$
of~\eqref{eq:partner-recurrence}. That the two coincide is the
\emph{sequence-level} form of Theorem~\ref{thm:sym2-template}, and it is not a
formality: the operator identity $\Lthree = P_2\cdot\Symtwo(\Ltwo)$ says that
$\Symtwo(\Ltwo)$ annihilates $f^2$, and passing from there to a statement
about coefficient sequences is exactly the step that a proof assistant does
not grant for free.

\begin{theorem}[Sequence-level symmetric square; status (K)]\label{thm:bridge}
Let $(a,b,c,d)\in\Z^4$ and let $s\colon\mathbb N\to\Z$ satisfy the Cooper
recurrence
\begin{equation}\label{eq:cooper-recurrence}
  (n+1)^3 s(n+1) = (2n+1)\bigl(an^2+an+b\bigr)s(n) - n\bigl(cn^2+d\bigr)s(n-1)
  \qquad(n\geq1),
\end{equation}
the coefficient form of $\Lthree(a,b,c,d)\sum s(n)z^n=0$, with $s(0)=1$,
$s(1)=b$. Let $(a_n)$ be
the partner sequence defined by~\eqref{eq:partner-recurrence} with $a_0=1$,
$a_1=\tfrac b2$, and $(g_n)$ the formal square root of $\sum s(n)z^n$. Then
$a_n = g_n$ for every $n$. Equivalently
\[
  \sum_{i=0}^{n} a_i\,a_{n-i} \;=\; s(n)\qquad(n\geq 0).
\]
\end{theorem}

\begin{proof}
Write $y_n=\sum_{i=0}^n a_ia_{n-i}$. Three steps.

\emph{(i) Symmetrization.} For any kernel $F$ and any sequence $a$,
$\sum_{i=0}^{m}F(i,m-i)\,a_ia_{m-i}
  = \sum_{i=0}^{m}\tfrac12\bigl(F(i,m-i)+F(m-i,i)\bigr)a_ia_{m-i}$,
by reflecting the summation index.

\emph{(ii) The convolution satisfies the order-3 recurrence.} Multiply
\eqref{eq:partner-recurrence} at index $t=n-i$ by $(6i+2t+4)\,a_i$ and sum
over $0\le i\le n$. The three resulting sums have kernels
$(6j+2k)k^2$, \ $(6j+2k+2)(2ak^2+ak+\tfrac b2)$ and
$(6j+2k+4)(ck^2+ck+\tfrac{c+d}4)$ on the antidiagonals $j+k=n+2,\,n+1,\,n$
respectively, and their symmetrizations are \emph{constant} on each
antidiagonal, equal to
\[
  (n+2)^3,\qquad (2n+3)\bigl(a(n+1)^2+a(n+1)+b\bigr),\qquad
  (n+1)\bigl(c(n+1)^2+d\bigr).
\]
These are precisely the coefficients of the Cooper recurrence. The boundary
terms $i=n+1,n+2$ cancel using $a_1=\tfrac b2a_0$. Hence $(y_n)$ satisfies the
Cooper recurrence.

\emph{(iii) Uniqueness.} $y_0=1=s(0)$ and $y_1=2a_0a_1=b=s(1)$, and the leading
coefficient $(n+2)^3$ never vanishes, so $y=s$. Finally $\sum a_nz^n$ and
$\sum g_nz^n$ are two power series over $\Q$ with constant term $1$ and equal
squares; since $\Q[[z]]$ is a domain and their sum is a unit, they are equal.

Formalized as \artifact{partner\_eq\_sqrt} in
\artifact{Agora/Sequences/SqrtBridgeGeneric.lean}, with step~(ii) as
\artifact{conv\_cooper\_of\_rec\_generic} over arbitrary
$(A,B,C,D)\in\Q^4$; the weight $6i+2t+4$ is the entire creative content, and
the kernel verifies the rest by \artifact{ring}. The instances are
\artifact{partner\_eq\_sqrt\_s7} and \artifact{partner\_eq\_sqrt\_s10}, which
discharge the recurrence hypothesis by the WZ certificates of
\S\ref{sec:formalization}. No instance is stated for $s_{18}$: the development
contains no closed form for $s_{18}$ to instantiate at, and we do not assume
one.
\end{proof}

\begin{corollary}[status (K)]\label{cor:partner-dyadic}
Under the hypotheses of Theorem~\ref{thm:bridge}, every partner coefficient
$a_n$ lies in $\Z[\tfrac12]$. In particular this holds for the $s_7$ and
$s_{10}$ partners at every index.
\end{corollary}

So powers of $2$ in the denominators are the \emph{forced} baseline for the
partner of any Cooper-template instance, now as a theorem rather than as an
observation: $s_{10}$ exhibits it (Theorem~\ref{thm:non-integral} shows the
prime $2$ genuinely occurs), and it is $s_7$ that is anomalous. For $s_7$ the
corollary excludes every odd prime from the denominators unconditionally; it
does \emph{not} give integrality, which remains the content of
Corollary~\ref{cor:s7-integral} and rests on the citation discussed there.

The 2-adic half can be isolated exactly.

\begin{proposition}[status (K)]\label{prop:mod4}
Let $s\colon\mathbb N\to\Z$ with $s(0)=1$ and $4\mid s(n)$ for all $n\geq1$. Then
every coefficient $g_n$, $n\geq1$, of the formal square root is an even integer.
Consequently, if $4\mid s_7(n)$ for all $n\geq1$, the $s_7$ partner is integral,
with no appeal to Theorem~\ref{thm:obrien}.
\end{proposition}

\begin{proof}
Strong induction: if $g_k\in2\Z$ for $1\le k<n$ then
$\sum_{k=1}^{n-1}g_kg_{n-k}\in4\Z$, so
$g_n=\tfrac12\bigl(s(n)-\sum g_kg_{n-k}\bigr)\in2\Z$. The second statement is
Theorem~\ref{thm:bridge}. Formalized as \artifact{sqrtSeq\_even\_of\_four\_dvd}
and \artifact{s7\_partner\_integral\_of\_congruence} in
\artifact{Agora/Sequences/SqrtIntegrality.lean}.
\end{proof}

\begin{remark}[PASS(200)]\label{rem:pass200}
The hypothesis $4\mid s_7(n)$ holds for $1\leq n\leq200$ in exact integer
arithmetic, the minimum 2-adic valuation over that range being exactly $2$; it
is kernel-checked for $n\leq6$. This is evidence, not proof, and
Corollary~\ref{cor:s7-integral} does not rely on it. If proved, it would replace
the development's only literature axiom by an elementary congruence on the
binomial sum defining $s_7$. The criterion is sufficient but not necessary
($(1+z)^2$ has an integral square root), so it is not a characterization; it
does fail for $s_{10}$ already at $n=1$ (\artifact{s10\_fails\_criterion}).
\end{remark}

\begin{remark}[A cautionary note on ``blocked'' goals]\label{rem:blocked}
For $s_7$, Theorem~\ref{thm:bridge} stood for some weeks as the development's
last open goal, with four recorded failed strategies and an internal ruling
that it was blocked on the absence, in the pinned library, of a power-series
square root and of a holonomic-sequence API. Both absences are real and neither
was relevant: the proof above uses only reindexing of finite sums and the fact
that $\Q[[z]]$ is a domain. We record this because the failure mode is general.
A ``blocked on the library'' verdict is a statement about a proof route; four
failed strategies are evidence about four strategies.
\end{remark}

\subsection{Non-integrality for \texorpdfstring{$s_{10}$}{s10} and
\texorpdfstring{$s_{18}$}{s18}}

\begin{theorem}[status (K)]\label{thm:non-integral}
The holomorphic solution of the $s_{10}$ partner has coefficients beginning
$1,\;1,\;\tfrac{17}{2},\;\tfrac{147}{2}$, and that of the $s_{18}$ partner
begins $1,\;3,\;\tfrac{45}{2},\;\tfrac{429}{2}$. In particular neither
sequence is integral.
\end{theorem}

\begin{proof}
Evaluate~\eqref{eq:partner-recurrence} at the respective parameters; the terms
listed are exact rationals. Integrality of the $n=2$ term would force
$17 = 2m$ resp. $45 = 2m$ for an integer $m$, which is impossible. The
formalization states the refutation in this arithmetic form rather than as
denominator bookkeeping: \artifact{s10\_partner\_two\_not\_integral},
\artifact{s18\_partner\_two\_not\_integral}, and the conclusions
\artifact{s10\_partner\_not\_integral}, \artifact{s18\_partner\_not\_integral}
in \artifact{Agora/Sequences/PartnerIntegrality.lean}.
\end{proof}

Note the asymmetry of difficulty: refuting integrality needs one witness,
establishing it needs all of them. The formalization keeps that asymmetry
visible in the type of its summary statement
(\artifact{cooper\_partner\_separation}), which asserts only the two negative
halves.

\subsection{Integrality for \texorpdfstring{$s_7$}{s7}}

For $(a,b,c,d) = (13,4,-27,3)$, recurrence~\eqref{eq:partner-recurrence}
becomes, after clearing,
\begin{equation}\label{eq:s7-partner-recurrence}
  (k+2)^2a_{k+2}
  = \bigl(26(k+1)^2 + 13(k+1) + 2\bigr)a_{k+1}
    + 3(3k+1)(3k+2)\,a_k,
  \qquad a_0=1,\ a_1=2 .
\end{equation}
Its first terms are
\[
  1,\;2,\;22,\;336,\;6006,\;117348,\;2428272,\;52303680,\;\dots
\]
which is OEIS A279619~\cite{OEIS}. Integrality of this sequence is a theorem
of the literature, and we state it as such.

\begin{theorem}[cited, O'Brien 2016]\label{thm:obrien}
\cite[Thm.~6.2, p.~47]{OBrien2016}. Any rational sequence satisfying
recurrence~\eqref{eq:s7-partner-recurrence} with the stated initial values is
integer-valued.
\end{theorem}

O'Brien's proof is an explicit coefficient-matching induction using the
$q$-expansions of two specific weight-1 level-7 modular objects; it is not a
formal consequence of any general integrality principle, and \S\ref{ssec:mechanism}
explains why. We do not reprove it.

\begin{corollary}[status (K), modulo Theorem~\ref{thm:obrien}]
\label{cor:s7-integral}
The holomorphic solution of the $s_7$ partner is integral, and its coefficient
sequence is A279619.
\end{corollary}

\begin{proof}
That the sequence defined by the generic
recurrence~\eqref{eq:partner-recurrence} at the $s_7$ parameters satisfies
exactly~\eqref{eq:s7-partner-recurrence}, with the same initial values, is a
mechanical identity; it is kernel-checked as
\artifact{partnerSeq\_s7\_recurrence}. Theorem~\ref{thm:obrien} then applies.
In the formalization the cited theorem is the registered axiom
\artifact{obrien2016\_theorem6\_2}, and the conclusion is
\artifact{s7\_partner\_integral}; see \S\ref{sec:formalization} for the axiom's
statement and its disclosure.
\end{proof}

\begin{remark}[PASS(7)]\label{rem:pass7}
The formalization independently retains a kernel-checked finite verification,
\artifact{s7\_partner\_integral\_pass7}: the coefficients are integers for
$n\leq7$, each obtained by evaluating the recurrence rather than by table
lookup, with the eight witnesses displayed above. This was the honest position
before the citation was traced to its correct source, and it is kept on the
record. It is evidence, not proof, of the general statement, and is reported
as \emph{PASS(7)} rather than as a theorem about all $n$
(\S\ref{sec:formalization}).
\end{remark}

\subsection{The mechanism: a normalized integral uniformizer}
\label{ssec:mechanism}

Why is $s_7$'s partner integral when the dyadic floor of
Proposition~\ref{prop:dyadic} permits denominators? The answer is not that
$\eta$-quotients have integer Fourier coefficients. It is a normalization
statement, and the distinction is not pedantic --- the object in question is a
series \emph{reversion}, and reversion does not preserve integrality in
general.

Following~\cite{OBrien2016}, set
\begin{equation}\label{eq:obrien-objects}
  z_7 = \sum_{m,n\in\Z}q^{\,m^2+mn+2n^2}, \qquad
  Z_7 = \frac{7P(q^7)-P(q)}{6}, \qquad
  X_7 = \frac{\eta_1^3\eta_7^3}{z_7^{\,3}},
\end{equation}
with $P$ the Ramanujan Eisenstein series $P(q) = 1-24\sum_{n\geq1}\sigma(n)q^n$
and $\eta_1^3\eta_7^3 = q\prod_{j\geq1}(1-q^j)^3(1-q^{7j})^3$. O'Brien's
Theorem~6.1 gives the expansion $z_7 = \sum_{n\geq0}c_7(n)X_7^{\,n}$, and the
$c_7(n)$ are the terms of~\eqref{eq:s7-partner-recurrence}.

\begin{proposition}[Reversion of a normalized integral series]
\label{prop:normalized-reversion}
Let $X \in q\,\Z[[q]]$ have leading coefficient exactly $1$, i.e.
$X = q + O(q^2)$, and let $Y \in \Z[[q]]$. Then the unique expansion
$Y = \sum_{n\geq0}c_nX^n$ has $c_n \in \Z$ for all $n$.
\end{proposition}

\begin{proof}
Determine the $c_n$ recursively: $c_0 = Y(0)$, and having fixed
$c_0,\dots,c_{n-1}$, the residual $Y - \sum_{k<n}c_kX^k$ lies in
$q^n\Z[[q]]$, while $X^n = q^n + O(q^{n+1})$; so $c_n$ is the coefficient of
$q^n$ in the residual divided by $1$, hence an integer. Integrality of the
leading coefficient of $X$ is exactly what makes that division trivial; a
leading coefficient of $2$, say, would reintroduce denominators of precisely
the kind Proposition~\ref{prop:dyadic} permits.
\end{proof}

\begin{proposition}[The mechanism, verified; status (E)]\label{prop:mechanism}
With $X_7$ as in~\eqref{eq:obrien-objects}, exact rational arithmetic on
$q$-expansions gives
\[
  X_7 \;=\; q - 9q^2 + 30q^3 - 15q^4 + \cdots \in q\,\Z[[q]],
\]
so $X_7$ is a normalized integral uniformizer; and expanding $z_7$ in powers
of $X_7$ reproduces $1,2,22,336,6006,117348,2428272$, i.e. A279619.
Consequently, by Proposition~\ref{prop:normalized-reversion}, integrality of
$c_7$ follows from~\eqref{eq:obrien-objects} together with O'Brien's
Theorem~6.1 --- but from those taken \emph{jointly}, and not from the
integrality of an $\eta$-quotient alone.
\end{proposition}

\begin{remark}[Resolving the source's equation (3.15); status (E)]
\label{rem:315}
The pinned copy of~\cite{OBrien2016} renders the denominator of its equation
(3.15) ambiguously between $z_7$ and $Z_7$ and drops the exponent, so which
object is meant is not determined by the text. Three readings were tested in
exact arithmetic:
\[
  \frac{\eta_1^3\eta_7^3}{z_7^{\,3}}, \qquad
  \frac{\eta_1^3\eta_7^3}{Z_7^{\,3}}, \qquad
  \frac{\eta_1^3\eta_7^3}{Z_7^{\,2}} .
\]
\emph{All three} are normalized integral uniformizers, and all three therefore
produce integral expansion coefficients --- which is exactly why integrality
alone cannot resolve the ambiguity. They produce, respectively,
\[
  1,2,22,336,6006,117348,2428272; \qquad
  1,2,34,828,23662,739984,24521448; \qquad
  1,2,26,476,10206,239280,5942232,
\]
so only the first reproduces A279619. It is also the only weight-consistent
reading: $\eta_1^3\eta_7^3$ has weight $3$ and $z_7$ has weight $1$, so
$z_7^{\,3}$ is what makes the quotient weight $0$, as a uniformizer must be.
The computation additionally verifies $z_7^{\,2} = Z_7$ exactly, confirming
that the two symbols denote genuinely different objects. The resolution is
recorded so that the ambiguity need not be re-litigated.
\end{remark}

\begin{remark}[Scope]\label{rem:integrality-scope}
Proposition~\ref{prop:mechanism} establishes nothing about $s_{10}$ or
$s_{18}$. Their partners are non-integral by Theorem~\ref{thm:non-integral},
and no modular parametrization by a normalized integral uniformizer is known
for them --- but absence of a known parametrization is not a proof that none
exists, and no such negative claim is made here. In particular no level is
assigned to $s_{18}$ anywhere in this paper.
\end{remark}

\subsection{The modular coordinate}

The parameter of the $s_7$ family is a modular function for the same group
that Remark~\ref{rem:signature} identified from the exponent data.

\begin{proposition}[status (E), with a cited group-theoretic input]
\label{prop:hauptmodul}
Let $t$ be the generating function of OEIS A279618~\cite{OEIS} and let
$u = q^{-1}\bigl(E(q)/E(q^7)\bigr)^4$, with $E(q)=\prod_{n\geq1}(1-q^n)$, be
the classical $\Gamma_0(7)$ Hauptmodul; put $v = 1/u \in q\,\Z[[q]]$. Then,
in exact rational arithmetic on $q$-expansions:
\begin{enumerate}
\item $t$ is \emph{not} a M\"obius function of $v$ (the linear system solved
  from the first 6 orders admits no solution verifying on later orders);
\item $t$ \emph{is} of the form $(\alpha v^2+\beta v+\gamma)/(\delta
  v^2+\varepsilon v+\zeta)$, with coefficients solved linearly from the first
  12 orders and then verified on held-out orders 13 and 29 that were not used
  in solving (29 orders available in total);
\item rewriting $t$ as a degree-2 rational function of $u$ and solving for the
  involution constant $\kappa$ with $R(\kappa/u)\equiv R(u)$ yields
  $\kappa = 49$.
\end{enumerate}
Items (i) and (ii) give $[\Q(u):\Q(t)] = 2$, so $\Q(t)$ is the fixed field of
an element normalizing $\Gamma_0(7)$. The normalizer of $\Gamma_0(N)$ in
$\mathrm{PSL}_2(\mathbb{R})$ for prime $N$ is the Fricke extension
$\Gamma_0(N)+$ \cite{AtkinLehner1970} --- this step is cited, not re-derived
--- so that element is the Fricke involution and $t$ is a $\Gamma_0(7)+$
Hauptmodul. Item (iii) exhibits the Fricke constant $\kappa = 7^2$ directly.
\end{proposition}

The constant $\kappa$ is computed from the fitted coefficients, never assumed;
the checker's verdict requires it to come out $49$ but does not supply that
value to the fit. Three negative controls accompany the computation: feeding
it $v$ itself (a genuine $\Gamma_0(7)$ Hauptmodul) must return the verdict
``not plus'', and does; corrupting one coefficient of A279618 must break the
degree-2 verification, and does; and testing against the wrong level,
$v_5 = q\bigl(E(q^5)/E(q)\bigr)^6$ for $\Gamma_0(5)$, must fail every fit, and
does.

\begin{remark}[Three independent indicators of the level]
\label{rem:level-indicators}
The number $7$ has now appeared from three computations that share no code:
the implied Fuchsian signature of the $s_7$ operators, matching
$\Gamma_0(7)+$ (Remark~\ref{rem:signature}, status (E)); the Hauptmodul and
Fricke constant $\kappa = 49$ of Proposition~\ref{prop:hauptmodul} (status
(E)); and the $\sqrt7$ appearing uninvited in the numerically continued
monodromy, together with $\det G = -14$ (Numerical Claim~\ref{nc:monodromy}
and Proposition~\ref{prop:lattice}, status (N)/(E)). What that convergence
does and does not license is the subject of \S\ref{sec:framework}.
\end{remark}

% ----------------------------------------------------------------------------
% Generated-by: Opus 5 (Stream 1 paper agent) | Verified-by: recurrence and
% values vs Agora/Sequences/PartnerIntegrality.lean (partnerPair, s7/s10/s18_
% partner_values, partnerSeq_s7_recurrence, s7_partner_integral_pass7) and
% Agora/Axioms/OBrien2016.lean; dyadic baseline vs Agora/Sequences/FormalSqrt.
% lean (sqrtSeq_dyadic, dyadic_baseline_control); X7/readings/c7 vs Stream 2
% data/certificates/S7_PARTNER_MODULAR.json; Hauptmodul orders, held-out orders,
% kappa and controls vs HAUPTMODUL_S7_GAMMA07PLUS.json and
% checkers/check_s7_hauptmodul_gamma07plus.py. Claims #7-#10 of paper/PLAN.md §3.
% | Reviewed-by: pending T0 (Xavier)
% ----------------------------------------------------------------------------

% CONDITIONAL SECTION — included under Option A of paper/PLAN.md §1.3 (scope
% ruling still open for T0). Under Option B this file is deleted wholesale,
% together with 06-lattice.tex; no other section depends on it except via the
% label sec:framework.
% ============================================================================
% 08-dolgachev-doran.tex — the M_n-polarized framework. Claims #18-#20 of
% PLAN.md §3. rho=19 / T=3 is presented CONDITIONALLY (PLAN.md §5 item 4); the
% identification of the computed lattice with T is conjectural throughout.
% ============================================================================
\section{Relation to the \texorpdfstring{$M_n$}{Mn}-polarized framework}
\label{sec:framework}

This section places the computations of \S\S\ref{sec:sym2}--\ref{sec:integrality}
next to a classical framework and states carefully what the juxtaposition does
and does not establish. Nothing in this section is claimed as a new theorem.

\subsection{The framework, as cited}

\begin{theorem}[cited, Dolgachev 1996]\label{thm:dolgachev}
\cite[\S7, p.~20 and Thm.~7.1]{Dolgachev1996}. For $M_n = \U\oplus E_8^2
\oplus\lat{-2n}$ one has $(M_n)^{\perp} = \U\oplus\lat{2n}$ in the K3 lattice,
and the coarse moduli space of $M_n$-polarized K3 surfaces is the Fricke
modular curve $\mathbb{H}/\Gamma_0(n)+$.
\end{theorem}

\begin{theorem}[cited, Doran 2000]\label{thm:doran}
\cite[Thm.~5.13]{Doran2000}. The Picard--Fuchs operator of an
$M_n$-polarized family of K3 surfaces is the symmetric square of a
second-order Fuchsian operator.
\end{theorem}

Both sources were fetched, hash-pinned, and read by the program; the parts of
Theorem~\ref{thm:dolgachev} used below were additionally re-derived
symbolically before use, with the period shape and the invariant bilinear form
recomputed from scratch rather than transcribed. We also record, because it
governs how much the coincidences below can be made to prove, that Doran
himself notes the classification of rank-19 lattice polarizations to be
incomplete~\cite[\S6]{Doran2000}.

\subsection{The logical direction}

A referee's first reaction to \S\ref{sec:lattice} should be that
Theorem~\ref{thm:dolgachev} already says $(M_7)^{\perp} = \U\oplus\lat{14}$
and Theorem~\ref{thm:doran} already says the Picard--Fuchs operator of such a
family is a symmetric square, so nothing has been learned. The reply is that
the direction of inference here is the reverse of the framework's. The
framework starts from a polarization hypothesis and deduces invariants. We
start from a concretely given operator family --- specified by four rational
numbers --- and \emph{compute} its invariants: the symmetric-square structure
(kernel-checked, and for the whole template rather than for one member), the
Fuchsian signature, the modular coordinate and its level, and the monodromy
lattice. Each of these is a fact about the operator, established without any
geometric hypothesis. The framework then supplies a reading under which they
all cohere, and the coherence is close: the same level $7$ appears from four
computations that share no code, and the monodromy lattice matches
$(M_7)^{\perp}$ on the nose, including the discriminant form.

That is convergent evidence, not a proof, and the incompleteness of the
rank-19 classification is precisely why it cannot be upgraded by citation.

\subsection{The evidence, assembled}

\begin{table}[htbp]
\centering
\begin{tabular}{@{}p{0.42\textwidth}ll@{}}
\toprule
\textbf{Computed fact about the $s_7$ operators} & \textbf{Status} &
\textbf{Framework counterpart} \\
\midrule
$\Lthree = P_2\cdot\Symtwo(\Ltwo)$, uniformly on the template
  (Thm.~\ref{thm:sym2-template}) & (K) & Thm.~\ref{thm:doran} \\
Almkvist--van Straten $W\equiv0$ (Thm.~\ref{thm:avs}) & (K) &
  self-adjointness of a $\Symtwo$ \\
$\Ltwo$, $\Lthree$ irreducible; $\Lthree$ minimal
  (\S\ref{sec:irreducibility}) & (E) & rank-3 sub-system \\
Implied Fuchsian signature: genus $0$, orders $(2,2,3)$, one cusp,
  area$/2\pi = \tfrac23$ (Rem.~\ref{rem:signature}) & (E) &
  signature of $\Gamma_0(7)+$ \\
Modular coordinate is a $\Gamma_0(7)+$ Hauptmodul, $\kappa = 49$
  (Prop.~\ref{prop:hauptmodul}) & (E) & $\mathbb{H}/\Gamma_0(7)+$ of
  Thm.~\ref{thm:dolgachev} \\
$\sqrt7$ in the elliptic monodromies (NC~\ref{nc:monodromy}) & (N) &
  level $7$ \\
Monodromy lattice $\U\oplus\lat{14}$, $\det = -14$, disc.\ $\Z/14$,
  $q = \tfrac1{14}$ (Props.~\ref{prop:lattice},~\ref{prop:usplit}) & (N)+(E) &
  $(M_7)^{\perp}$ of Thm.~\ref{thm:dolgachev} \\
$s_{10}$ control gives $\U\oplus\lat{20}$ (NC~\ref{nc:s10-control}) & (N)+(E) &
  level $10$, not level $7$ \\
Orthogonal complement of $G$ in $\Lambda$ is
  $\U\oplus E_8(-1)^2\oplus\lat{-14}$ (Prop.~\ref{prop:g0complement}) & (N)+(E) &
  exactly $M_7$, mutual complement of Thm.~\ref{thm:dolgachev} \\
\bottomrule
\end{tabular}
\medskip
\caption{Computed invariants of the $s_7$ operator family against the
$M_n$-polarized framework. Status labels as in \S\ref{sec:intro}.}
\label{tab:evidence}
\end{table}

\subsection{A conditional rank statement}

We now record the rank statement that the minimality result of
\S\ref{sec:irreducibility} yields \emph{if} two external inputs are granted.
It is stated as a conditional proposition, with its hypotheses displayed, and
it is used nowhere else in this paper.

\begin{proposition}[conditional; status (C)+(E) under (H1)--(H2)]
\label{prop:ranks}
Assume:
\begin{itemize}
\item[\textbf{(H1)}] the $s_7$ family is realized by the explicit projective
  K3 model of Almkvist and van Straten~\cite{AlmkvistvanStraten2021}, and
  $\Lthree$ is the Picard--Fuchs operator of the period of its holomorphic
  $2$-form; and
\item[\textbf{(H2)}] for the very general member $X$ of the family,
  $\operatorname{rank}T(X)$ equals the rank of the $\Q$-sub-variation of Hodge
  structure generated by that period under Gauss--Manin.
\end{itemize}
Then $\operatorname{rank}T(X) = 3$ and $\rho(X) = 19$.
\end{proposition}

\begin{proof}
By Corollary~\ref{cor:minimal}, the sub-variation generated by the period has
rank $3$, that being the order of its minimal annihilator. (H2) transfers this
to $\operatorname{rank}T(X)$, and $\rho(X) = 22 - \operatorname{rank}T(X) = 19$.
\end{proof}

\begin{remark}[On (H2), and on the caveat that must travel with $\rho = 19$]
\label{rem:H2}
(H2) is standard Hodge theory and we cite rather than re-derive it. It is the
conjunction of two inclusions. First, for very general $z$ the N\'eron--Severi
lattice is locally constant, the Noether--Lefschetz locus being a countable
union of proper subvarieties, so $T = \NS^{\perp}$ is a Gauss--Manin-flat
sub-Hodge-structure containing the holomorphic $2$-form; the sub-variation
generated by that form is therefore contained in $T$. Second, $T(X)\otimes\Q$
is by definition the minimal primitive sub-Hodge-structure of $H^2(X,\Q)$
whose complexification contains $H^{2,0}$, and the saturation of the
sub-variation is such a structure; so the reverse inclusion holds
\cite[ch.~3 \S3.1, ch.~17 \S17.1]{Huybrechts2016}. Simplicity of the
transcendental Hodge structure is~\cite[Thm.~1.6(a)]{Zarhin1983}.

The caveat is essential and is not a formality: $\rho = 19$ is the value for
the \emph{very general} member. It jumps to $20$ on a countable dense subset
of the base. Any use of Proposition~\ref{prop:ranks} that drops the phrase
``very general'' is a misuse of it.
\end{remark}

\subsection{An independent cross-check: the orthogonal complement (G0)}
\label{subsec:g0}

A separate, purely lattice-theoretic computation supplies additional support
for the coherence recorded in Table~\ref{tab:evidence}, and it is reported
here with exactly the same conditionality as \S\ref{sec:lattice}: it takes the
recognized matrices of Numerical Claim~\ref{nc:monodromy} as input, is
otherwise exact, and is dated 2026-07-28.

\begin{proposition}[The orthogonal complement of $G$ in the K3 lattice;
status (E), given Numerical Claim~\ref{nc:monodromy}]\label{prop:g0complement}
Let $G$ be the Gram matrix of Proposition~\ref{prop:lattice}(iii), embedded
primitively in the K3 lattice $\Lambda = \U^3\oplus E_8(-1)^2$ via the
hyperbolic-plane witness $f\mapsto e_{U_1},\ e\mapsto f_{U_1},\ w\mapsto
e_{U_2}+7f_{U_2}$ (a norm-$14$ vector in a hyperbolic plane is forced to have
this shape; there is no other primitive choice). Its orthogonal complement is
isometric over $\Z$ to
\[
  \U \oplus E_8(-1) \oplus E_8(-1) \oplus \lat{-14},
\]
of rank $19$, signature $(1,18)$, cyclic discriminant group of order $14$
(discriminant-form value $q = \tfrac{27}{14} \bmod 2\Z$ on a generator), and
existence and uniqueness of the primitive embedding are certified rather than
assumed, via the Nikulin/Eichler bound $\operatorname{rank}\Lambda -
\operatorname{rank}G = 19 \ge \ell(A_G) + 2 = 3$.
\end{proposition}

Two remarks fix what this proposition does and does not add to
Proposition~\ref{prop:ranks}.

\emph{What it does not add.} Arithmetically, rank $19$ here is forced by
$\operatorname{rank}\Lambda - \operatorname{rank}G = 22 - 3$; it carries no
information about $\rho$ beyond what Proposition~\ref{prop:ranks} already
gives, and it is \emph{not} presented as a second derivation of $\rho = 19$.
Nor does it discharge hypotheses (H1)/(H2) of Proposition~\ref{prop:ranks}, or
either gap of Conjecture~\ref{conj:T} below: it is a further exact computation
downstream of the same numerically-recognized input (Numerical
Claim~\ref{nc:monodromy}), not an independent route to the identification of
the transcendental lattice.

\emph{What it does add.} The proposition computes the \emph{explicit isometry
type} of the complement, not merely its rank. Comparing it to
Theorem~\ref{thm:dolgachev}: under the standard convention that $E_8^2$
there denotes $E_8(-1)\oplus E_8(-1)$ (as fixed by $\Lambda$ itself, above),
$\U\oplus E_8(-1)^2\oplus\lat{-14}$ is exactly $M_7$. On the computation, $G$
and $M_7$ are thus mutual orthogonal complements inside $\Lambda$ --- the same
coherence already recorded for $T$ and $(M_7)^{\perp}$ in
Table~\ref{tab:evidence}, now confirmed from the other side, together with a
verified (not assumed) uniqueness-of-embedding bound.

On independence, precisely: this computation has been carried out three
times --- two in-house exact-symbolic implementations sharing no code (a
constructive embedding witness, and an independent genus-uniqueness argument),
agreeing on the Gram matrix bit-for-bit, and one further zero-shot
re-derivation by a separate reviewing model, which built its own $E_8$ Cartan
matrix from a different Dynkin-node labeling and independently recomputed the
discriminant value; all three agree exactly
\cite{G0NSGenus2026,G0DecisionLog2026}. This is reported as corroboration of
the computation, not as a fourth independent geometric derivation of $\rho=19$:
the third check was supplied the specific embedding formula displayed above
rather than asked to discover it, and, like the first two, it starts from the
same Numerical Claim~\ref{nc:monodromy} input.

\subsection{The lattice arithmetic, kernel-checked}\label{subsec:mn-lean}

The identification of the computed lattice with $\U\oplus\lat{14}$ is status
(E) and is not re-derived in Lean. What \emph{is} kernel-checked
(\artifact{Agora/Geometry/MnLattice.lean}) is the lattice arithmetic that the
identification is then fed into, uniformly in $N$. It is built on the lattice
interface of an independent Lean development of the K3 and Narain
lattices~\cite{LeanMaster2026}, which supplies the Gram-matrix type, the
hyperbolic plane, and the signature of $\Lambda_{K3}$; nothing in the present
subsection depends on that development's physical content.

\begin{proposition}[status (K)]\label{prop:mn-lean}
Let $N\in\Z$ and let $T_N$ be the Gram matrix of $\U\oplus\lat{2N}$ in a basis
$(e,f,w)$ with $e^2=f^2=0$, $e\cdot f=1$, $w^2=2N$.
\begin{enumerate}
\item $T_N$ is symmetric and even, with $\det T_N=-2N$; for $N=7$ it is not
  unimodular.
\item \emph{(Glue.)} In a single hyperbolic plane, $e+Nf$ and $e-Nf$ are
  orthogonal of norms $2N$ and $-2N$, and span a sublattice of index $2N$;
  $e+Nf$ is primitive. For $N=7$ this is the witness $w\mapsto e+7f$ of
  Proposition~\ref{prop:g0complement}.
\item \emph{(Swap.)} The exchange $\sigma\colon e\leftrightarrow f$ is an
  isometry of $T_N$, an involution, of determinant $-1$.
\item \emph{(Swap $=$ Fricke.)} Over any field, the vector
  $\omega(\tau)=e-N\tau^2f+\tau w$ is isotropic, and for $N,\tau\neq0$
  \[
    \sigma\,\omega(\tau) \;=\; (-N\tau^2)\;\omega\!\Bigl(-\frac{1}{N\tau}\Bigr).
  \]
  If $N\tau^2=-1$ then $\sigma\,\omega(\tau)=\omega(\tau)$.
\end{enumerate}
\end{proposition}

Item~(iv) is the statement, at the level of the period line, that the lattice
involution $\sigma$ induces the Fricke involution $\tau\mapsto-1/(N\tau)$ on
the tube-domain parameter of the $M_N$-polarized period domain
\cite{Dolgachev1996}; the formalization includes a control refuting the
variant with the opposite sign, so the statement is not an artifact of the
normalization. The signatures $(2,1)$ and $(1,18)$ assigned to
$\U\oplus\lat{2N}$ and $M_N$ add to the $(3,19)$ of $\Lambda_{K3}$
(\artifact{sig\_fills\_K3}); as in the cited development, these are asserted
pairs, not extracted from definiteness proofs, and the rank $19$ is forced by
$22-3$.

One structural observation, stated with its limits. The $2\times2$ block of
$\sigma$ on the $\U$ summand is the Gram matrix of $\U$ itself, which the cited
development proves equal to the Narain form of a circle compactification, where
the same exchange (of momentum and winding) is the generator of T-duality. So a
single integer matrix is, on one lattice, the Fricke involution of the $M_N$
period domain and, on another, the T-duality generator. This is a fact about a
lattice isometry. We do not assert, and nothing here supports, a physical
identification of the two.

\subsection{The modular group as a group of lattice isometries}\label{subsec:modular-action}

Proposition~\ref{prop:mn-lean}(iii)--(iv) exhibited one isometry of
$\U\oplus\lat{2N}$ acting on the period as the Fricke involution. That was a
sample of a general structure, which we now state in the form in which it is
kernel-checked (\artifact{Agora/Geometry/ModularAction.lean}). Write
$\omega(\tau) = e - N\tau^2 f + \tau w$ as before.

\begin{proposition}[status (K)]\label{prop:rho}
Over any commutative ring, define
\[
  \rho(a,b,c,d) \;=\;
  \begin{pmatrix}
    d^2 & -Nc^2 & 2Ncd\\
    -Nb^2 & a^2 & -2Nab\\
    bd & -ac & ad+Nbc
  \end{pmatrix},
  \qquad
  \rho_W(a,b,c,d) \;=\;
  \begin{pmatrix}
    Nd^2 & -c^2 & 2Ncd\\
    -b^2 & Na^2 & -2Nab\\
    bd & -ac & Nad+bc
  \end{pmatrix}.
\]
Then:
\begin{enumerate}
\item $\rho(a,b,c,d)$ is an isometry of $\U\oplus\lat{2N}$ whenever
  $(ad-Nbc)^2=1$, and $\rho_W(a,b,c,d)$ is one whenever $(Nad-bc)^2=1$.
\item $\rho$ is multiplicative: $\rho(g)\rho(g') = \rho(gg')$, the product being
  that of $\Gamma_0(N)$ written on the entries.
\item \emph{(Automorphy.)} As an identity of polynomials --- no determinant
  hypothesis, no invertibility ---
  \[
    \rho(a,b,c,d)\,\omega(\tau)
      \;=\; \bigl((Nc\tau+d)^2,\; -N(a\tau+b)^2,\; (a\tau+b)(Nc\tau+d)\bigr),
  \]
  which over a field is $(Nc\tau+d)^2\,\omega\!\bigl(\tfrac{a\tau+b}{Nc\tau+d}\bigr)$.
  Similarly $\rho_W$ induces $\tau\mapsto\frac{Na\tau+b}{N(c\tau+d)}$.
\end{enumerate}
\end{proposition}

So $\rho$ realizes $\Gamma_0(N)$, and $\rho_W$ the Atkin--Lehner elements
$W=\left(\begin{smallmatrix}Na&b\\Nc&Nd\end{smallmatrix}\right)/\sqrt N$, inside
$O(\U\oplus\lat{2N})$, acting on the period line with the weight-2 automorphy
factor; the $\sqrt N$ cancels in the symmetric square, which is why
Atkin--Lehner involutions are lattice isometries at all. This is the
constructive half of Dolgachev's $O^{+}(\U\oplus\lat{2n})/\pm1\cong\Gamma_0(n)^{+}$
\cite{Dolgachev1996}; the converse is quoted, not proved. The Fricke involution
is $\rho_W(0,-1,1,0)$, which is $-\sigma$ for the $\sigma$ of
Proposition~\ref{prop:mn-lean}, so that proposition is the specialization of this
one. A control confirms the determinant hypothesis in (i) is load-bearing: at
$(a,b,c,d)=(1,0,0,2)$, $N=7$, the matrix is not an isometry.

\subsection{The singular points of the operator are self-dual points}\label{subsec:selfdual}

We can now say what the two finite singular points of $\Lthree$ are. Recall from
\S\ref{sec:operators} that the leading coefficient of the $s_7$ partner operator
is $P_2 = 1-26z-27z^2$, whose roots are $z=-1$ and $z=\tfrac1{27}$; an earlier
reading of these as Kodaira degenerations was withdrawn, and they were
subsequently identified as order-2 elliptic points of $X_0(7)^{+}$.

\begin{proposition}[status (K)]\label{prop:selfdual}
Let $z(h) = h/(1+13h+49h^2)$. Then, over any field:
\begin{enumerate}
\item $z\bigl(\tfrac1{49h}\bigr) = z(h)$, so $z$ is invariant under the Fricke
  involution of the $h$-line;
\item over $\Q$, the fixed points of $h\mapsto \tfrac1{49h}$ are exactly
  $h=\pm\tfrac17$, and $z(\tfrac17)=\tfrac1{27}$, $z(-\tfrac17)=-1$;
\item $\bigl(1-26z(h)-27z(h)^2\bigr)\,(1+13h+49h^2)^2 = (1-49h^2)^2$.
\end{enumerate}
Moreover $\langle e-f,\,\omega(\tau)\rangle = -(N\tau^2+1)$, so the locus
$N\tau^2=-1$ fixed by the Fricke involution is exactly the wall of the
$(-2)$-root $e-f$; and for $N=7$ the conjugate Atkin--Lehner element
$\rho_W(1,4,-2,-1)$ is an involution whose negative is the reflection in the
$(-2)$-root $(-2,4,1)$ of $\U\oplus\lat{14}$.
\end{proposition}

Item (iii) is the sharp statement: pulled back to the $h$-line, the leading
coefficient of the operator becomes a \emph{perfect square} vanishing exactly at
the two Fricke fixed points. The singular locus of $\Ltwo,\Lthree$ is the
ramification locus of the quotient by the involution, and both of its points are
walls of $(-2)$-roots of the lattice --- one for the Fricke involution, one for
its conjugate. We regard this as the correct replacement for the withdrawn
Kodaira reading: it is an exact statement, it is kernel-checked, and it explains
why there are exactly two such points and why they are the ones they are.

\begin{remark}[What the identification costs]\label{rem:selfdual-status}
Proposition~\ref{prop:selfdual} is a theorem about the rational map $z(h)$ and
about the lattice. Reading it as a statement about the $s_7$ \emph{family}
requires two further inputs, neither kernel-checked. First, that
$h=\bigl(\eta(7\tau)/\eta(\tau)\bigr)^4$ and $z=z(h)$ pull the family back to the
$\tau$-line: we verify the equivalent $q$-series identity
$\sum_n s_7(n)z(q)^n = \bigl(\sum_{a,b}q^{a^2+ab+2b^2}\bigr)^2$ exactly to
$O(q^{40})$, status (E), with a control in which replacing $13$ by $12$ makes the
identity fail. Second, that $h(-1/(7\tau)) = 1/(49h(\tau))$, which is the
$\eta$~transformation law, status (L); numerically $h(i/\sqrt7)=\tfrac17$ and
$h\bigl((-1+i/\sqrt7)/2\bigr)=-\tfrac17$ to 40 digits, and the Fricke relation
holds to $10^{-42}$ at a generic point. The script is
\artifact{scripts/check\_selfdual\_points\_s7.py}.
\end{remark}

\begin{remark}[No physical claim]\label{rem:selfdual-nophysics}
In the string-theory literature the wall of a $(-2)$-root is associated with
enhanced gauge symmetry, and the fixed locus of T-duality on a circle is the
self-dual radius; the $2\times2$ block of $\sigma$ on the $\U$ summand is
literally the Narain Gram matrix of a circle, on which the same exchange is the
T-duality generator. We record this because it is what makes the structure
interesting, and we claim nothing from it: no compactification is constructed
here, and a coincidence of lattice isometries is not a physical identification.
\end{remark}

\subsection{What is conjectural}

\begin{conjecture}[Identification of the computed lattice]\label{conj:T}
The joint monodromy-invariant lattice computed in \S\ref{sec:lattice} is the
transcendental lattice of the very general member of the $s_7$ family; that
is, $T \cong \U\oplus\lat{14}$, and the family is $M_7$-polarized.
\end{conjecture}

We are explicit about the two gaps between Proposition~\ref{prop:lattice} and
Conjecture~\ref{conj:T}, because within the computation itself the
identification is pinned only up to two enumerable ambiguities, and only one
of them is closed by computation.

\begin{enumerate}
\item \emph{Even invariant overlattices.} A priori the transcendental lattice
  could be an even overlattice of the computed orbit lattice, with the same
  rational form. These were enumerated: there are none
  (Proposition~\ref{prop:lattice}(v)). This gap is closed, conditionally on
  Numerical Claim~\ref{nc:monodromy}.
\item \emph{Overall integral rescaling.} The transcendental lattice could
  carry $\lambda\cdot G$ for an integer $\lambda>1$. This branch is excluded
  by the shape $\U\oplus\lat{2n}$ that the framework predicts
  (Theorem~\ref{thm:dolgachev}) --- that is, by a reading, not by a
  computation. It is the honest residual gap, and we do not paper over it: the
  argument that rules it out presupposes the conclusion it is used to support.
\end{enumerate}

To these must be added the fact that Numerical Claim~\ref{nc:monodromy} is
itself numerics-backed, and hypothesis (H1) of
Proposition~\ref{prop:ranks}, which is a citation about a geometric model and
is not verified anywhere in this paper. Proposition~\ref{prop:g0complement}
(\S\ref{subsec:g0}) closes neither gap: it computes the orthogonal complement
of the same conjectured lattice, and so is exactly as conditional as the
lattice itself.

\subsection{What is not claimed}

We close by listing, explicitly, statements a reader might expect to find here
and will not.

\begin{itemize}
\item No degeneration or fibre-type reading is made of the singular points
  $z=-1$ and $z=\tfrac1{27}$. Remark~\ref{rem:signature} identifies them as
  order-2 elliptic points of the quotient curve; that is the only reading
  taken, and no elliptic fibration is invoked anywhere.
\item No claim is made that the $s_{10}$ or $s_{18}$ partners are modular,
  and no level is assigned to $s_{18}$ (Remark~\ref{rem:integrality-scope}).
\item Theorem~\ref{thm:sym2-template} is not claimed to imply
  Theorem~\ref{thm:doran} or to be implied by it: the template theorem
  exhibits a partner from four rational parameters with no geometric
  hypothesis, whereas Doran's theorem asserts existence under a polarization
  hypothesis. Neither subsumes the other.
\item Proposition~\ref{prop:ranks} is not asserted unconditionally, and
  Conjecture~\ref{conj:T} is not asserted at all.
\item Proposition~\ref{prop:g0complement} (\S\ref{subsec:g0}) is not claimed
  as evidence independent of \S\ref{sec:lattice}: it shares Numerical
  Claim~\ref{nc:monodromy} with every other statement in this section, and its
  three-way agreement is a corroboration of one exact computation, not three
  independent geometric derivations of $\rho=19$.
\item \textbf{[STUB, 2026-07-30]} No claim of the form ``$T(\text{$s_{10}$}) \cong
  \U\oplus\lat{20}$'' is made anywhere in this paper, and none is available to
  make: the only extant occurrence of $\U\oplus\lat{20}$ in the program is
  Numerical Claim~\ref{nc:s10-control} of \S\ref{sec:lattice}, a
  \emph{discriminating control} run on the identical pipeline to show
  level-sensitivity against the $s_7$ result (Remark~\ref{rem:lattice-scope}
  states explicitly that even the $s_7$ lattice's identification with a
  transcendental lattice is not established there, let alone $s_{10}$'s). It
  is not a certified, independently-verified monodromy-lattice derivation for
  $s_{10}$ analogous to Proposition~\ref{prop:usplit}/\ref{rem:two-verifications}
  for $s_7$: no second implementation, no negative-control suite, and no
  $s_{10}$-specific G0-style orthogonal-complement check exist for it in this
  repository or in Stream~2's epistemic ledger (\texttt{PREDICTION.md},
  \texttt{PROOF\_STATUS.txt}, \texttt{K3\_SELECTION\_REPORT.md} --- checked
  2026-07-30, no hit beyond the control note itself and an identical
  regression-control comment in Stream 2's G0 checker). Any inbound brief
  presenting $s_{10}$'s $\U\oplus\lat{20}$ as a certified, stand-alone result
  should be treated as stale and returned for provenance, per the program's
  standing verification rule. \emph{[Update 2026-08-01: Stream~2 has since
  promoted the fiber-level Nikulin-complement computation for $s_{10}$ to a
  titled draft certificate (\texttt{G0\_NS\_genus\_cooper\_s10.json}, Tier~B,
  pending review), with its own negative-control suite. That certificate
  takes the $s_{10}$ transcendental-lattice identification as an \emph{input}
  and states explicitly that the reviewed monodromy-lattice derivation this
  item says is missing remains missing; the stance of this item --- no
  $s_{10}$ lattice claim is made in this paper --- is unchanged.]}
\end{itemize}

\subsection{Open questions}

\begin{enumerate}
\item Make Numerical Claim~\ref{nc:monodromy} exact. The natural route is
  exact hypergeometric connection formulae for $\Ltwo$, which would upgrade
  Propositions~\ref{prop:lattice} and~\ref{prop:usplit} to unconditional
  statements about the $s_7$ family and remove the only numerical input in
  this paper.
\item Close the rescaling ambiguity of Conjecture~\ref{conj:T} by a
  computation rather than by the framework's shape --- for instance by
  exhibiting a $\lambda>1$ obstruction intrinsic to the monodromy
  representation.
\item Prove Theorem~\ref{thm:obrien} in a proof assistant, discharging the one
  literature axiom of the formal development (\S\ref{sec:formalization}). This
  requires $q$-expansion machinery for weight-1 level-7 modular objects that
  the pinned library does not have.
\item Determine whether the $s_{18}$ partner exhibited in
  Corollary~\ref{cor:sporadic-partners} --- which appears not to have been
  written down before --- has any modular interpretation, or whether its
  non-integrality is an obstruction to one.
\item Establish the genus statement for the abstract lattice
  $\U\oplus\lat{14}$ independently. It is not needed for
  Proposition~\ref{prop:usplit}, which is constructive and therefore stronger
  than a genus argument for the specific lattice at hand, but it would be the
  natural framing for a classification statement.
\end{enumerate}

\subsection*{A further construction (placeholder, 2026-07-29)}

A further section, extending the lattice-theoretic results above toward an
explicit higher-dimensional construction, is planned but not yet drafted: an
in-progress feasibility check on the relevant base spaces has not concluded,
and drafting ahead of it risks stating a result that the check may
contradict. This placeholder will be replaced once that check resolves; no
content, method, or conclusion of the pending section is anticipated here.

% ----------------------------------------------------------------------------
% Generated-by: Opus 5 (Stream 1 paper agent), G0/G1 additions 2026-07-29 by
% Sonnet 5 (WP-P1) | Verified-by: framework statements vs paper/references.bib
% entries Dolgachev1996 / Doran2000 as recorded in PLAN.md row #19; rho=19 /
% T=3 conditionality and the (H1)/(H2) split vs Stream 2
% data/certificates/L3_IRREDUCIBLE.json step5_hodge_identification +
% conditional_on_step5; overlattice and rescaling ambiguities vs
% C2_cooper_s7_v4.json and briefs/STREAM2_U1_EXECUTION_2026_07_27.md §4.
% Claims #18-#20 of paper/PLAN.md §3; presented conditionally per PLAN.md §5.4.
% G0 (Prop.~g0complement, \S8.4): S2 briefs/G0_NS_GENUS_RESULT_2026_07_28.md
% (commits 9a386d9, 15f16c5) + decision log
% briefs/WP_S2G_X4_EXHIBITION_PLAN_2026_07_27.md §8 (commit 256017d, LIVE
% promotion + cross-model verification entry); independence caveats
% transcribed from the G0 brief's own "Two honest limits" paragraph, not
% invented here. G1 placeholder: PLAN.md scope ruling (pure mathematics, zero
% physics) kept in mind -- deliberately undescribed and unmotivated per WP-P1
% instructions; whether it belongs in this manuscript at all is flagged as a
% scope question for T0 in briefs/P1_DRAFT_STATUS_2026_07_29.md.
% U+<20> stub added 2026-07-30 by Sonnet 5 (T0 Directive 3): searched this
% repo (paper/PLAN.md row 17, 06-lattice.tex nc:s10-control) and Stream 2
% (~/SocrateAI-Scientific-Agora-K3-DarkMatter) PREDICTION.md, PROOF_STATUS.txt,
% K3_SELECTION_REPORT.md for a certified cooper_s10 monodromy-lattice
% derivation; none found -- only the discriminating-control number (already
% correctly tiered NC/(N)+(E) in this manuscript) and an identical
% regression-control comment in Stream 2's checkers/check_NS_genus_G0.py /
% data/certificates/G0_NS_genus_cooper_s7.json. Per Stream 2 Standing Rule 4,
% the directive's implied "certified cooper_s10 section" is not executed;
% this explicit stub is inserted instead. Deviation-with-cause recorded here
% and in the commit message.
% | Reviewed-by: pending T0 (Xavier)
% ----------------------------------------------------------------------------

% ============================================================================
% 09-formalization.tex — the Lean 4 development. Toolchain facts checked against
% lean-toolchain (v4.34.0-rc2) and lake-manifest.json (mathlib tag v4.34.0-rc2,
% rev 85e3a25e...). Updated 2026-09-20 with the toolchain migration; the prior
% pin was v4.32.0 / mathlib 3dffaf2f. The migration required no source change
% and left the axiom inventory and the native_decide disclosure unchanged --
% each re-verified at the new pin, not carried over.
% See briefs/MEMO_LEAN_4_34_MIGRATION_2026_09_20.md.
% FURTHER UPDATE 2026-09-20: open_goal_partner_eq_sqrt_s7 was CLOSED
% (Agora/Sequences/SqrtBridge.lean); the development now has ZERO sorry. The
% "blocked-on-mathlib" ruling on it was wrong -- see
% briefs/BRIDGE_GOAL_CLOSED_2026_09_20.md. Sorry-count claims updated throughout.
% Axiom inventory checked against AXIOMS.md and Agora/Axioms/*.lean.
% ============================================================================
\section{The formal development}\label{sec:formalization}

The statements labeled \textbf{(K)} in this paper are theorems of a Lean~4
development which we describe here in enough detail to be audited. We also
describe what the kernel did \emph{not} check, since a formalization paper
that reports only its successes is of no use to a referee.

\subsection{Toolchain and layout}

The development builds with Lean~4 at toolchain \texttt{v4.34.0-rc2}
\cite{deMouraUllrich2021} against \texttt{mathlib4}~\cite{mathlib4,mathlibCPP2020}
pinned at tag \texttt{v4.34.0-rc2}, commit
\texttt{85e3a25e006c35636f0e53b0e9296caca2685bc0}. (The development was
migrated to this pin from Lean~\texttt{v4.32.0} / \texttt{mathlib4}
\texttt{3dffaf2f} on 2026-09-20; no source file required any change, and every
statement in the development is unchanged, as certified by a
statement-hash lock taken before the migration and re-checked after it.) It is
organized so that each epistemic category has exactly one place to live:

\begin{center}
\begin{tabular}{@{}ll@{}}
\toprule
Directory & Contents and rule \\
\midrule
\artifact{Agora/} & The mathematics. No \artifact{sorry} anywhere. \\
\artifact{Agora/Axioms/} & The only permitted location for an
  \artifact{axiom} declaration. \\
\artifact{OpenGoals/} & The only permitted location for a \artifact{sorry};
  each is a named goal. \\
\artifact{Tests/} & Golden numeric tests against independently computed
  values. \\
\bottomrule
\end{tabular}
\end{center}

% ⚠️ T0 CORRECTION 2026-09-20 (Opus 5) — AMENDED ON T0 INSTRUCTION, same date.
% The previous wording claimed both placement rules were "enforced mechanically
% by a repository hook" and referred to a "no-sorry check". Verified false this
% session; corrected below. Record of what was wrong:
%   * "enforced mechanically by a repository hook" — the only mechanism is
%     .claude/hooks/lean_guard.sh, a Claude Code PostToolUse hook. It fires ONLY
%     when Claude edits a .lean file in a session loading .claude/settings.json.
%     A hand edit, another tool, or any `git commit` bypasses it entirely.
%   * "an axiom ... is rejected" — TRUE (the hook exits 2 on that path).
%   * "the no-sorry check excludes OpenGoals/ and nothing else" — there is NO
%     blocking no-sorry check. The hook only warns and exits 0, and the repo has
%     no .github/ directory, so there is no CI at all. (CLAUDE.md rule 3 made the
%     same "CI-enforced" claim; it is corrected on this branch.)
% The paper must not claim an enforcement mechanism stronger than the one that
% exists; that is the claim a referee is most entitled to check.
The two placement rules are enforced differently, and we state the difference
because a formalization paper that overstates its own safeguards forfeits the
audit it is asking for. The \artifact{axiom} rule is enforced mechanically: an
editor hook rejects an \artifact{axiom} declaration introduced outside
\artifact{Agora/Axioms/}. The \artifact{sorry} rule is \emph{not} mechanically
enforced --- it is maintained by review, and checked by the audit commands of
\S\ref{subsec:not-checked}, which report every occurrence and its location; at
the time of writing they report none at all. We can now be sharper about
\emph{why} the build is not the gate here, because it was tested rather than
assumed: appending a \artifact{sorry}-carrying theorem to a module and
rebuilding exits \emph{zero}, reporting success with only a
\texttt{declaration uses `sorry`} warning. What does catch it is the axiom
audit, via \artifact{sorryAx}. A reader auditing this development should
therefore run the audit and not infer a clean development from a clean build. This matters because the failure mode
being guarded against is not dishonesty but drift --- an assumption introduced
locally during a proof attempt and never surfaced.

\subsection{The results, and how they are stated}

The kernel-checked content of this paper is:

\begin{center}
\begin{tabular}{@{}lll@{}}
\toprule
Statement & Lean name & File \\
\midrule
Thm.~\ref{thm:sym2-template}, $\thetaop^3$ &
  \artifact{partner\_res3} & \artifact{Sequences/PartnerOperators.lean} \\
Thm.~\ref{thm:sym2-template}, $\thetaop^2$ &
  \artifact{partner\_res2} & \artifact{"} \\
Thm.~\ref{thm:sym2-template}, $\thetaop^1$ &
  \artifact{partner\_res1} & \artifact{"} \\
Thm.~\ref{thm:sym2-template}, $\thetaop^0$ &
  \artifact{partner\_res0} & \artifact{"} \\
Lemma~\ref{lem:magic} &
  \artifact{partner\_magic} & \artifact{"} \\
Lemma~\ref{lem:c2-collapse} &
  \artifact{cooperC2\_eq\_thetaC\_cooperC3} & \artifact{"} \\
Cor.~\ref{cor:sporadic-partners} &
  \artifact{s7\_P2\_eq}, \dots, \artifact{s18\_P0\_eq} & \artifact{"} \\
Thm.~\ref{thm:avs} &
  \artifact{P\_cleared\_eq\_zero} & \artifact{Swampland/SymSquareC3b.lean} \\
Rem.~\ref{rem:golden-sym2} &
  \artifact{symSquare\_golden\_harmonic/\_exp} & \artifact{SymSquare.lean} \\
Prop.~\ref{prop:dyadic} &
  \artifact{sqrtSeq\_dyadic} & \artifact{Sequences/FormalSqrt.lean} \\
Thm.~\ref{thm:non-integral} &
  \artifact{s10\_partner\_not\_integral}, \artifact{s18\_\dots} &
  \artifact{Sequences/PartnerIntegrality.lean} \\
Cor.~\ref{cor:s7-integral}, mechanical half &
  \artifact{partnerSeq\_s7\_recurrence} & \artifact{"} \\
Cor.~\ref{cor:s7-integral} &
  \artifact{s7\_partner\_integral} & \artifact{"} \\
Rem.~\ref{rem:pass7} &
  \artifact{s7\_partner\_integral\_pass7} & \artifact{"} \\
Thm.~\ref{thm:bridge}, step (ii) &
  \artifact{conv\_cooper\_of\_rec\_generic} &
  \artifact{Sequences/SqrtBridgeGeneric.lean} \\
Thm.~\ref{thm:bridge} &
  \artifact{partner\_eq\_sqrt}, \artifact{\_s7}, \artifact{\_s10} & \artifact{"} \\
Cor.~\ref{cor:partner-dyadic} &
  \artifact{partner\_dyadic}, \artifact{s10\_partner\_dyadic} & \artifact{"} \\
Prop.~\ref{prop:mn-lean} &
  \artifact{glue\_congruence}, \artifact{swap\_is\_fricke}, \dots &
  \artifact{Geometry/MnLattice.lean} \\
Thm.~\ref{thm:mod4} &
  \artifact{four\_dvd\_s7},
  \artifact{s7\_partner\_integral\_axiom\_free} &
  \artifact{Sequences/S7Mod4.lean} \\
Prop.~\ref{prop:embedding} &
  \artifact{B\_pullback}, \artifact{C\_pullback},
  \artifact{B\_orthogonal\_C} & \artifact{Geometry/Embedding.lean} \\
Rems.~\ref{rem:two-lattices}, \ref{rem:signature-tierA} &
  \artifact{no\_isometry\_G0N\_TN}, \artifact{sym2\_isometry\_general},
  \artifact{TN\_diagonalises} &
  \artifact{Geometry/SymSquareForms.lean}, \artifact{MnLattice.lean} \\
Prop.~\ref{prop:rho} &
  \artifact{rho\_isometry}, \artifact{rho\_mul},
  \artifact{rho\_mulVec\_period} & \artifact{Geometry/ModularAction.lean} \\
Prop.~\ref{prop:selfdual} &
  \artifact{zOf\_fricke}, \artifact{s7\_P2\_discriminant},
  \artifact{W7conj\_eq\_neg\_reflection} &
  \artifact{Geometry/SelfDual.lean}, \artifact{"} \\
\bottomrule
\end{tabular}
\end{center}

Three design decisions shape how these are stated.

\emph{$\thetaop$-form, not monic $D$-form.} Normalizing an order-3 operator to
be monic in $D$ introduces genuine rational-function coefficients, and the
rational-function type in the pinned library carries no derivative API. Rather
than improvise one, every statement above lives in \artifact{Polynomial}~$\Q$:
no division, no nonvanishing side conditions, and every proof obligation a
polynomial identity. This is why Theorem~\ref{thm:sym2-template} is stated as
four coefficient identities plus the collapse lemma rather than as an operator
equation --- the factorization in $\Q(z)[\thetaop]$ is then derived on paper,
as in \S\ref{sec:sym2}, from kernel-checked polynomial facts.

\emph{Concrete identities, not existence claims.} The relation
``$\Lthree$ is the symmetric square of $\Ltwo$'' is encoded as an equality
fixing \emph{every} coefficient of $\Lthree$ as an explicit polynomial function
of those of $\Ltwo$. It is falsifiable per candidate. The development had an
earlier generation of statements of the form ``there exists a polynomial of
degree 3'', true of any cubic and therefore content-free; those were deleted
and the episode is recorded in the repository's escalation log. The lesson
generalizes: a formalized statement is worth exactly the content of its
statement, and no amount of kernel checking supplies content the statement
lacks.

\emph{Derivatives are computed, not transcribed.} Every occurrence of a
derivative in Theorem~\ref{thm:avs} is an application of the library's
\artifact{Polynomial.derivative}, so the kernel performs each differentiation
itself. Encoding a hand-computed derivative as an opaque definition would let
a transcription error pass \artifact{ring} undetected.

\subsection{The axiom inventory, in full}\label{subsec:axioms}

The development registers exactly one load-bearing literature axiom.

\begin{center}
\begin{minipage}{0.93\textwidth}
\small
\begin{verbatim}
axiom obrien2016_theorem6_2 (c : N -> Q) (h0 : c 0 = 1) (h1 : c 1 = 2)
    (hrec : forall k : N,
      ((k : Q) + 2)^2 * c (k + 2) =
        (26 * ((k : Q) + 1)^2 + 13 * ((k : Q) + 1) + 2) * c (k + 1)
          + 3 * (3*(k : Q) + 1) * (3*(k : Q) + 2) * c k) :
    forall n, exists m : Z, c n = (m : Q)
\end{verbatim}
\end{minipage}
\end{center}

\noindent
This is Theorem~\ref{thm:obrien} verbatim, re-indexed by $n = k+1$ to avoid
truncated subtraction on $\mathbb{N}$ and with the source's boundary
computation replaced by the two initial values as hypotheses. It is used
exactly once, to conclude \artifact{s7\_partner\_integral} from
\artifact{partnerSeq\_s7\_recurrence}. It is registered in the repository's
axiom register with its provenance, and the source was read in full before
being cited --- which mattered: an earlier attribution in the program named
the wrong theorem of the same source and a mechanism (``$\eta$-quotients have
integer coefficients, hence a corollary'') that does not in fact establish
integrality, since the object is a series reversion. See
Proposition~\ref{prop:normalized-reversion} for what the correct mechanism is.

It is not provable in the pinned library: O'Brien's argument needs the
$q$-expansions of two specific weight-1 level-7 modular objects, obtained by
classical modular-forms machinery that \texttt{mathlib} at the pinned commit
does not contain. Discharging it is open question~3 of
\S\ref{sec:framework}.

For completeness: the register carries one further axiom declaration,
\artifact{pipeline\_upper\_bound}, belonging to a part of the repository
unrelated to the mathematics of this paper. It is disclosed in the register as
vacuously true --- its statement is an existential satisfied by the witness
$1$ and encodes no data --- and no file used in this paper imports it. We
mention it because an axiom inventory that reports only the axioms one wishes
to discuss is not an inventory.

\subsection{Golden validation and PASS(\texorpdfstring{$N$}{N}) discipline}

Two devices keep the development from certifying its own mistakes.

\emph{Golden tests} check a formula against an answer obtained by other means.
The symmetric-square formula~\eqref{eq:sym2-formula} is validated against two
examples whose answers are derivable by inspection
(Remark~\ref{rem:golden-sym2}); no candidate-level symmetric-square theorem is
accepted until those are green, because the formula is the single point of
failure for every such theorem. The sequence values in \artifact{Tests/} are
checked against reference values computed independently in exact integer
arithmetic outside Lean, for $n \leq 19$. Separately, the recurrence-defined
partner sequence is checked against the formal square root of the bulk
sequence, on both the integral side ($s_7$: $2, 22, 336$) and the
non-integral side ($s_{10}$: $1, \tfrac{17}{2}$) --- so the $\tfrac{17}{2}$
witnessing Theorem~\ref{thm:non-integral} is not an artifact of one encoding.

\emph{PASS($N$)} is the reporting convention for a kernel-checked finite
range. \artifact{s7\_partner\_integral\_pass7} verifies integrality for
$n\leq7$ and is cited as PASS(7), never as ``the sequence is integral'' --- a
distinction that is not pedantry here, since establishing integrality needs
every term while refuting it needs one. The convention survived the arrival of
the axiom-backed general theorem: the finite statement was kept rather than
deleted, and the summary theorem \artifact{cooper\_partner\_separation}
deliberately asserts only the two \emph{negative} halves of the separation, so
that no downstream file can cite ``only $s_7$ is integral'' as a single
uniformly-established fact.

\subsection{The three-strikes escalation discipline}
\label{subsec:three-strikes}

A lemma is not worked indefinitely. The development's house rule --- three
named strategies attempted and recorded, then escalate or name an open goal,
never grind unboundedly --- governs every stalled proof attempt, and it is
mechanical rather than a stylistic preference: an escalation path routes a
lemma from the working tier, through one tier up after three failed
attempts \emph{or} a verifier-uncheckable obstruction, to the top after a
further two failed attempts \emph{or} a genuine design ambiguity, with
de-escalation available once the top tier decomposes the problem.

The development's last open goal,
\artifact{open\_goal\_partner\_eq\_sqrt\_s7} (that the recurrence-defined
partner sequence equals the formal square root of the bulk $s_7$ series,
solution-level transport of Theorem~\ref{thm:sym2-template}), carries its
three-strikes record verbatim in \artifact{OpenGoals/PartnerIntegrality.lean}:
(1) direct strong induction on $n$, which fails because the inductive step
needs a convolution identity equivalent to the goal itself; (2) simultaneous
induction on the equality together with the convolution identity, which
reduces to a WZ-style certificate identity not derivable without an
operator-action-on-sequences bridge the pinned library lacks; (3) a Mathlib
search (\artifact{exact?}/\artifact{apply?}), which returns nothing, since no
\artifact{PowerSeries.sqrt} or holonomic-sequence API exists at the pinned
commit. All three are recorded as failures of route, not failures of effort.
The same file also records a fourth, post-ruling attempt (2026-07-26): the
underlying convolution identity was re-verified as an exact-arithmetic
identity to $n=38$ by two independent encodings --- evidence, not proof, that
the goal is true --- while a computer-algebra certificate search
(\artifact{ore\_algebra} symmetric power) produced an operator that
demonstrably does \emph{not} annihilate the convolution sequence, a mismatch
recorded in the file as unresolved rather than explained away. The goal was
carried as a named \artifact{sorry} in \artifact{OpenGoals/}, per the
placement rule of \S\ref{sec:formalization}, rather than being silently
dropped or its statement weakened.

That record is the reason we can report the sequel honestly. On 2026-09-20 the
goal was \emph{closed}, by a fifth route that none of the four recorded
attempts had considered: a convolution sum may be symmetrized in its two
indices, after which the solution-level transport is a single linear
combination against the order-2 recurrence. The ``blocked-on-library'' verdict
that attempt (3) had supported was therefore wrong --- and wrong for the second
time on this same pair of goals, the sibling
\artifact{open\_goal\_partner\_integral\_s7} having been closed earlier by a
citation the same verdict had ruled out. We draw the obvious conclusion rather
than a flattering one: the discipline's value is that it \emph{records
routes and preserves the statement}, not that its verdicts are reliable. Four
recorded failures are evidence about four routes. Because the statement had
never been weakened to accommodate them, the later proof closed the original
goal verbatim --- checked here by a statement lock against the untouched
declaration, together with a negative control confirming that the lock can
fail.

The discipline is not merely a stopping rule: the sibling goal
\artifact{open\_goal\_partner\_integral\_s7} was opened under an identical
three-strikes record (induction on the recurrence, a $p$-adic-valuation
strengthening, and a Mathlib search --- each with its own recorded failure
mode) and was later \textbf{closed} --- not by any of the three grind
strategies succeeding, but by locating the applicable literature theorem
(O'Brien's Theorem~6.2, cited as the disclosed axiom
\artifact{obrien2016\_theorem6\_2} of \S\ref{sec:formalization} above, not
re-proved) after escalation. The record of
the failed attempts was kept rather than deleted once the citation route
closed the goal, on the view that a documented dead end is itself
information for the next person who reaches the same lemma.

\subsection{What the kernel did not check}\label{subsec:not-checked}

\begin{itemize}
\item \emph{Theorem~\ref{thm:obrien}}, as discussed. Kernel-checked:
  that our sequence satisfies its hypotheses.
\item \emph{That Gorodetsky's operator is the right operator for the Cooper
  sequences.} Our $\theta$-coefficients \artifact{cooperC3}--\artifact{cooperC0}
  encode eq.~(1.7) of \cite{Gorodetsky2023}, and every symmetric-square identity
  of \S\ref{sec:sym2} is a statement about that encoding. The \emph{transcription}
  has been verified: expanding
  $\theta^3 - z(2\theta+1)(a\theta^2 + a\theta + b) + z^2(c(\theta+1)^3 + d(\theta+1))$
  and collecting powers of $\theta$ reproduces all four encoded coefficients,
  checked twice by independent routes with a negative control. What is
  \emph{not} re-derived here is that (1.7) is itself the correct operator for
  these sequences; that is Gorodetsky's result, cited. The distinction matters:
  a sign or commutator slip in the transcription would have left every identity
  in \S\ref{sec:sym2} true of an operator that is not Cooper's.
\item \emph{Everything with status \textbf{(E)} or \textbf{(N)}}:
  \S\ref{sec:irreducibility}, \S\ref{sec:lattice}, and the computational parts
  of \S\ref{sec:integrality} are Python and computer-algebra computations, not
  Lean. They are exact and reproducible (\S\ref{sec:reproducibility}) but they
  are not kernel-checked, and the paper never labels them as though they were.
\item \emph{The golden tests in} \artifact{Tests/} use
  \artifact{native\_decide} for terms up to $n=19$, whose binomial
  coefficients make kernel evaluation impractical. \artifact{native\_decide}
  trusts the compiler in addition to the kernel; both affected theorems carry
  an explicit trust annotation. No result in \S\S\ref{sec:sym2}--\ref{sec:integrality}
  depends on them.
\item \emph{No open goal remains.} The last one --- the equality, for all $n$,
  of the recurrence-defined partner sequence with the formal square root of the
  bulk sequence (\artifact{open\_goal\_partner\_eq\_sqrt\_s7}) --- was
  \textbf{closed on 2026-09-20}, and the development now contains no
  \artifact{sorry} at all. It is the solution-level transport of the operator
  identity that Theorem~\ref{thm:sym2-template} establishes at operator level.
  The proof (\artifact{Agora/Sequences/SqrtBridge.lean}) depends only on Lean's
  three standard axioms and does not use the citation axiom of
  \S\ref{subsec:axioms} below.

  We record why it had been open, because the reason is instructive and
  cautionary for this style of work. The goal had been ruled
  ``blocked-on-\texttt{mathlib}'' on the grounds that the library provides
  neither a formal-power-series square root nor a holonomic-sequence API. Both
  observations are correct and neither is relevant: the proof uses only
  \artifact{PowerSeries.mk}, \artifact{coeff\_mul}, antidiagonal reindexing,
  \artifact{isUnit\_iff\_constantCoeff}, the power-series domain instance, and
  \artifact{Finset.sum\_range\_reflect}, all of which were present in the
  library version pinned at the time the goal was declared blocked. The four
  recorded failed strategies were never refuted; they were circumvented. The
  step that had been missing is that a convolution sum may be symmetrized in
  its two indices, after which the solution-level transport --- which the
  earlier analysis had taken to be a creative-telescoping problem --- reduces to
  a single linear combination against the order-2 recurrence. A
  ``blocked-on-library'' verdict is a statement about a proof route, not about a
  goal; $n$ failed strategies are evidence about those $n$ strategies.
\end{itemize}

\subsection{An honest assessment of the formalization}

The proofs above are, as proofs, shallow: after the right statements are set
up, almost all of them are closed by \artifact{ring}. We do not claim
otherwise, and we do not think it detracts from the contribution. What is
difficult in this development is not discharging the goals but arriving at
statements that are simultaneously (i) true uniformly in the template
parameters rather than per candidate, (ii) division-free, so that they can be
stated at all in the available library, (iii) non-vacuous, in the sense that
they pin coefficients rather than assert existence, and (iv) honest about
which of them rests on a citation. The partner
of~\eqref{eq:generic-partner} being \emph{determined} by the template rather
than guessed, the collapse $\thetaop(P_2) = 2P_1$ making the identity
polynomial, and the axiom quarantine are the three places where that work
went. To our knowledge, no kernel-checked proof of a symmetric-square
structure theorem for Picard--Fuchs-type operators has previously been given
in any proof assistant.

\subsubsection*{A failure mode the gates do not catch}

One further disclosure belongs in an honest assessment, because it bears on how
much a reader should infer from the fact that this development passes its
checks. A systematic audit comparing each declaration's \emph{name} against
what its \emph{statement} actually mentions found nine instances of a single
defect: a theorem that is true, compiles, and passes the build, the
\artifact{sorry} grep, the axiom audit \emph{and} the statement lock, while
proving less than its name says --- the claim living in the identifier and the
docstring rather than in the statement. One example, since the shape is easier
to recognise than to describe: a lemma named for the primitivity of $e + Nf$ in
a hyperbolic plane had the statement $\mathrm{IsCoprime}\,(1 : \mathbb{Z})\,N$,
which is true for every $N$ in every commutative ring and mentions no vector,
no basis and no lattice.

Two things are worth saying about this. First, \textbf{none of the nine was
found by any gate}; all were found by reading statements against names, prompted
by an exchange with an independent formalization effort. A reader should weight
``passes every check'' accordingly --- the checks are necessary and they are not
sufficient, and this is a sharper statement of the vacuity caveat already made
in \S\ref{subsec:not-checked}. Second, and more usefully: \textbf{every one of
the nine occurred in the repository's legacy physics-facing modules, and none in
the arithmetic and lattice development that this paper reports}. Those modules
have since been separated into a quarantine directory that nothing in the core
imports; they are retained rather than deleted, and none of them is cited
anywhere in this paper. The results of \S\S\ref{sec:sym2}--\ref{sec:integrality}
and \S\ref{sec:framework} were subjected to the same audit and are
unaffected.

% ----------------------------------------------------------------------------
% Generated-by: Opus 5 (Stream 1 paper agent); \S9.4 (three-strikes discipline)
% added 2026-07-30 by Sonnet 5 (T0 Directive 3) | Verified-by: toolchain vs
% lean-toolchain (leanprover/lean4:v4.34.0-rc2) and lake-manifest.json (mathlib
% tag v4.34.0-rc2, rev 85e3a25e006c35636f0e53b0e9296caca2685bc0 -- re-verified
% at the new pin 2026-09-20 by Opus 5; axiom inventory, native_decide
% disclosure and sorry location all re-checked, all unchanged); axiom text transcribed from
% Agora/Axioms/OBrien2016.lean; second axiom + hook rules from AXIOMS.md;
% native_decide disclosure from Tests/CooperSequences.lean header; remaining
% sorry located by grep (OpenGoals/PartnerIntegrality.lean:201 only);
% three-strikes record transcribed verbatim from
% OpenGoals/PartnerIntegrality.lean (both goals, incl. the CLOSED one's
% original record) and the escalation-path wording cross-checked against
% CLAUDE.md rule 5 and EXECUTION_PLAN.md §"Escalation path".
% | Reviewed-by: pending T0 (Xavier)
% ----------------------------------------------------------------------------

% ============================================================================
% 10-reproducibility.tex — artifact index and negative-control philosophy.
% Repository URLs deliberately not typed: T0 supplies them at submission
% (paper/PLAN.md §4, "not in scope of this scaffold").
% ============================================================================
\section{Reproducibility}\label{sec:reproducibility}

Every number in this paper is produced by a script or a proof in one of two
repositories, referred to below as \textbf{R1} (the formal development, Lean~4)
and \textbf{R2} (the exact and numerical computations, Python with
\texttt{sympy} and \texttt{mpmath}). Repository locations are given with the
published version. Nothing in this paper is a constant typed from memory; the
program's standing rule is that a number may enter prose only by being read
from an artifact that a reader can regenerate.

\subsection{Artifact index}

\begin{center}
\small
\begin{tabular}{@{}llll@{}}
\toprule
Result & Status & Artifact & Repo \\
\midrule
Thms.~\ref{thm:sym2-template}, \ref{lem:magic}, Cor.~\ref{cor:sporadic-partners}
  & (K) & \artifact{Agora/Sequences/PartnerOperators.lean} & R1 \\
Thm.~\ref{thm:avs}
  & (K) & \artifact{Agora/Swampland/SymSquareC3b.lean} & R1 \\
Rem.~\ref{rem:golden-sym2}
  & (K) & \artifact{Agora/SymSquare.lean} \S3 & R1 \\
Prop.~\ref{prop:dyadic}
  & (K) & \artifact{Agora/Sequences/FormalSqrt.lean} & R1 \\
Thm.~\ref{thm:non-integral}, Cor.~\ref{cor:s7-integral}, Rem.~\ref{rem:pass7}
  & (K) & \artifact{Agora/Sequences/PartnerIntegrality.lean} & R1 \\
Thm.~\ref{thm:obrien} (as axiom)
  & (C) & \artifact{Agora/Axioms/OBrien2016.lean}, \artifact{AXIOMS.md} & R1 \\
Props.~\ref{prop:L3-scheme}, \ref{prop:L2-scheme}, Rem.~\ref{rem:signature}
  & (E) & \artifact{check\_L3\_riemann\_scheme.py} $\to$
          \artifact{L3\_RIEMANN\_SCHEME.json} & R2 \\
\S\ref{sec:irreducibility} in full
  & (E) & \artifact{check\_L3\_irreducible\_minimal.py} $\to$
          \artifact{L3\_IRREDUCIBLE.json} & R2 \\
\S\ref{sec:irreducibility} controls
  & (E) & \artifact{test\_L3\_irreducible\_minimal\_controls.py} & R2 \\
Props.~\ref{prop:mechanism}, Rem.~\ref{rem:315}
  & (E) & \artifact{check\_s7\_partner\_integrality\_modular.py} $\to$
          \artifact{S7\_PARTNER\_MODULAR.json} & R2 \\
Prop.~\ref{prop:hauptmodul}
  & (E) & \artifact{check\_s7\_hauptmodul\_gamma07plus.py} $\to$
          \artifact{HAUPTMODUL\_S7\_GAMMA07PLUS.json} & R2 \\
NC~\ref{nc:monodromy}, Prop.~\ref{prop:lattice}, NC~\ref{nc:s10-control}
  & (N)/(E) & \artifact{check\_U1\_lattice.py} $\to$
          \artifact{C2\_cooper\_s7\_v4.json} & R2 \\
\S\ref{sec:lattice} controls
  & (E) & \artifact{test\_U1\_controls.py} & R2 \\
Prop.~\ref{prop:usplit} (2nd implementation)
  & (E) & \artifact{checkers/check\_U1\_splitting\_independent.py} & R1 \\
Prop.~\ref{prop:usplit} controls
  & (E) & \artifact{checkers/test\_U1\_splitting\_independent\_controls.py} & R1 \\
Prop.~\ref{prop:g0complement}
  & (N)/(E) & \artifact{check\_NS\_genus\_G0.py} & R2 \\
Prop.~\ref{prop:g0complement} controls
  & (E) & \artifact{test\_NS\_genus\_G0\_controls.py} (5/5) & R2 \\
\bottomrule
\end{tabular}
\end{center}

The lattice certificate is content-addressed: the version underlying
\S\ref{sec:lattice} has
\[
  \text{SHA256} =
  \mathtt{036cd895db892aee9802514ec18668535f8896cee53cf6123533af9844387c3e},
\]
and the independent verification of Proposition~\ref{prop:usplit} was run
against that file and against its immediate predecessor, whose derived block
is byte-identical, with the same result. Literature sources used in this paper
were fetched and hash-pinned before being read; the hashes are recorded in the
repositories' provenance files. Proposition~\ref{prop:g0complement}'s
certificate and its two internal reports are likewise cited by name and commit
hash rather than folded silently into the surrounding prose
\cite{G0NSGenus2026,G0DecisionLog2026}.

\subsection{Running the artifacts}

R1 builds with \artifact{lake build}; the golden tests build with it. R2's
checkers are standalone scripts, each taking an optional \artifact{-{}-json}
argument that writes its certificate, so a checker can be re-run and its output
diffed against the committed certificate:
\begin{center}\small
\artifact{python3 checkers/check\_L3\_irreducible\_minimal.py -{}-json /tmp/out.json}
\end{center}
The independent splitting check takes the certificate path as an argument
rather than embedding it, and reads the target level from the certificate at
runtime:
\begin{center}\small
\artifact{python3 checkers/check\_U1\_splitting\_independent.py -{}-cert <path>}
\end{center}
The control suites are runnable directly or under \texttt{pytest}.

\subsection{Negative controls}

The program's operating rule is that \emph{a check which cannot fail is not a
check}. Every computation above ships with controls designed to make it fail,
and the controls are run, not merely written. The rule has a specific origin:
an earlier computation in the program passed its own test because its
statistic had been clamped below its own threshold, and the resulting claim had
to be retracted. The controls exist so that this cannot recur silently.

Concretely, the controls used in this paper are:

\begin{itemize}
\item \emph{Known-negative inputs.} The irreducibility checker is fed an
  operator with a hyperexponential solution by construction and one with
  distinct indicial roots at the origin; it must reject both, and does
  (\S\ref{sec:irreducibility}).
\item \emph{Cross-checker agreement.} The $\Ltwo$ exponents computed by one
  checker must reproduce, under pairwise summation, the $\Lthree$ Riemann
  scheme computed independently by another --- at all four singular points, for
  both families.
\item \emph{Machinery controls.} The numerical continuation must recover a
  matrix known analytically before it is used on unknown loops
  (Remark~\ref{rem:machinery-control}).
\item \emph{Corruption controls.} Perturbing a recognized matrix must break a
  downstream exact gate; three distinct corruptions are exercised, failing at
  the involution gate, at the infinity-product gate, and by collapsing the
  invariant-form solution space to dimension $0$ (\S\ref{sec:lattice}).
  Corrupting a coefficient of the modular series must break the degree-2
  verification (\S\ref{sec:integrality}).
\item \emph{Discrimination controls.} The identical lattice pipeline run on the
  $s_{10}$ family must produce a different, internally consistent answer, and
  does (NC~\ref{nc:s10-control}); the modular fit tested against the wrong
  level must fail every fit, and does; the genuine $s_7$ Gram matrix tested
  against $\U\oplus\lat{20}$ must fail, and does.
\item \emph{Vacuity controls.} A statement whose content is not obvious is
  accompanied by a check that it is not trivially true: the dyadic baseline of
  Proposition~\ref{prop:dyadic} is paired with a proof that
  $\tfrac13\notin\Z[\tfrac12]$.
\item \emph{Positive controls after corruption.} Every suite re-runs the
  unmodified input last, confirming that the scrambling hooks are surgical and
  leave no state behind.
\end{itemize}

Two further disciplines are worth stating because they constrain what the
artifacts are permitted to emit. First, target values are assembled at runtime
from data read out of the certificate under test, never hardcoded outside a
labeled control --- so a checker cannot accidentally test a constant against
itself. Second, checkers that cannot derive a quantity emit \texttt{null} for
it rather than a plausible value: the irreducibility certificate reports no
rank at all, precisely because the step that would license one is a citation
and not a computation (\S\ref{sec:framework}).

% ----------------------------------------------------------------------------
% Generated-by: Opus 5 (Stream 1 paper agent), G0 rows added 2026-07-29 by
% Sonnet 5 (WP-P1) | Verified-by: file names and certificate names checked
% against the two repositories' checkers/ and data/certificates/ directories;
% SHA256 from briefs/STREAM1_U1_INDEPENDENT_VERIFICATION_2026_07_27.md;
% control inventory from the control scripts themselves. G0 artifact names
% (check_NS_genus_G0.py, test_NS_genus_G0_controls.py, 5/5) from S2
% briefs/G0_NS_GENUS_RESULT_2026_07_28.md "Coordinator independent
% verification" paragraph. | Reviewed-by: pending T0 (Xavier)
% ----------------------------------------------------------------------------


\bibliographystyle{amsplain}
\bibliography{references}

\end{document}
% ---------------------------------------------------------------------------
% Generated-by: Fable 5 (Stream 1 paper agent) | Verified-by: pdflatex build,
% claims traced per paper/PLAN.md #3 | Reviewed-by: pending T0 (Xavier)
% ---------------------------------------------------------------------------
